%% 第13回編集委員会（物理第1巻）
%% p1-10.tex  力-10　板がたわむ：はりの曲げとたわみ（組 P1-6）
%% 数値はすべて Python で計算した（scratchpad/phys1/力-10/calc10.py・te10.py）。g＝9.8 m/s²。
%% 記号（計画書 C-1）：断面二次モーメント I、断面係数 Z、せん断力 Q、曲げモーメント M、分布荷重 w（N/m）、集中荷重 F、
%%   たわみ y（位置ごと）・δ（最大）、x は左の支点から右へ、y は下向きを正。曲率半径 ρ（力-1 ⑤）。
%% 出典：
%%   鋼のヤング率 205 GPa・すぎ 約7 GPa：巻末の表（ph1-back.tex）の値を使った。すぎの曲げヤング係数のばらつき
%%     （奈良県産スギ平角材 200本で平均 8.9 GPa、6.1〜12.4 GPa）：奈良県森林技術センター 資料
%%     https://www.pref.nara.jp/secure/238892/02magekyoudo.compressed.pdf
%%   SGP 50A（外径 60.5 mm・厚さ 3.8 mm・5.31 kg/m）、SGP 25A（34.0 mm・3.2 mm）：JIS G 3452（wermac.org の表で確認）
%%   床のはりのたわみの目安：平成12年建設省告示第1459号（変形増大係数2をかけたたわみがスパンの1/250以下）、
%%     木造のスパン表では床ばり 1/300 の例（三重県・群馬県のスパン表の解説）
%%   H形鋼 200×100×5.5×8：JIS G 3192 の標準寸法（I は角の丸みを無視して計算した値。表の値より少し小さい）
%%   2×4 材の断面 38 mm×89 mm：JAS（枠組壁工法構造用製材）の寸法型式 204
%%   ウォルフの法則（J. Wolff, Das Gesetz der Transformation der Knochen, 1892）
%%   ガリレオ『新科学対話』（1638）の片持ちはり、クーロン（1773）、ナビエ（1826）：Timoshenko, History of Strength of Materials (1953)
\makeatletter
% 小見出し（それだけで1行）
\newcommand\pxkhead[1]{\par\addvspace{4pt}\Needspace{6\baselineskip}\noindent\textbf{#1}\par\nobreak\noindent\ignorespaces}
% 単位の中の点（N·m など）
\newcommand\pxkq[2]{\qty[inter-unit-product=\ensuremath{\cdot}]{#1}{#2}}
\newcommand\pxku[1]{\unit[inter-unit-product=\ensuremath{\cdot}]{#1}}
% 支点（ピン）  \pxkpin{(x,y)}   はりの下の面が (x,y)
\newcommand\pxkpin[1]{\begin{scope}[shift={#1}]
  \filldraw[fill=bkLine!70, draw=bkInk, line width=0.6pt] (0,0) -- (-2.3,-3.8) -- (2.3,-3.8) -- cycle;
  \draw[bkInk, line width=0.6pt] (-3.4,-3.8) -- (3.4,-3.8);
  \foreach \i in {-3.0,-1.6,-0.2,1.2,2.6} {\draw[bkGray, line width=0.4pt] (\i,-3.8) -- ++(-1.1,-1.2);}
\end{scope}}
% 支点（ローラー）  \pxkroller{(x,y)}
\newcommand\pxkroller[1]{\begin{scope}[shift={#1}]
  \filldraw[fill=bkLine!70, draw=bkInk, line width=0.6pt] (0,0) -- (-2.3,-3.0) -- (2.3,-3.0) -- cycle;
  \draw[bkInk, line width=0.5pt] (-1.3,-3.6) circle (0.6) (0,-3.6) circle (0.6) (1.3,-3.6) circle (0.6);
  \draw[bkInk, line width=0.6pt] (-3.4,-4.2) -- (3.4,-4.2);
  \foreach \i in {-3.0,-1.6,-0.2,1.2,2.6} {\draw[bkGray, line width=0.4pt] (\i,-4.2) -- ++(-1.1,-1.2);}
\end{scope}}
% かべ（片持ちはりの根もと）  \pxkwall{(x,y)}{高さ}   (x,y) はかべの右の面のまん中
\newcommand\pxkwall[2]{\begin{scope}[shift={#1}]
  \fill[bkLine!60] (-3.5,-#2/2) rectangle (0,#2/2);
  \draw[bkInk, line width=0.8pt] (0,-#2/2) -- (0,#2/2);
  \pgfmathsetmacro\pxka{-#2/2}\pgfmathsetmacro\pxkb{-#2/2+2}\pgfmathsetmacro\pxkc{#2/2}\foreach \yy in {\pxka,\pxkb,...,\pxkc} {\draw[bkGray, line width=0.4pt] (0,\yy) -- ++(-2.2,-1.6);}
\end{scope}}
% 下向きの集中荷重の矢印  \pxkload{(x,y)}{長さ}{ラベル}   (x,y) は矢の先
\newcommand\pxkload[3]{\draw[bk arrow, bkOrange, line width=1.3pt] ($#1+(0,#2)$) -- #1;
  \node[anchor=south, font=\sffamily\footnotesize, text=bkOrange!70!black, inner sep=1pt] at ($#1+(0,#2+0.3)$) {#3};}
% 等分布荷重（小さな矢の列）  \pxkdist{x1}{x2}{はりの上の面の y}{ラベル}
\newcommand\pxkdist[4]{\fill[bkOrange!15] (#1,#3+1.0) rectangle (#2,#3+5);
  \draw[bkOrange, line width=0.6pt] (#1,#3+5) -- (#2,#3+5);
  \foreach \xx in {0,0.125,...,1.0} {\draw[-{Stealth[length=2.2pt]}, bkOrange, line width=0.6pt] ({#1+\xx*(#2-#1)},#3+5) -- ({#1+\xx*(#2-#1)},#3+0.4);}
  \node[anchor=south, font=\sffamily\footnotesize, text=bkOrange!70!black, inner sep=1pt] at ({(#1+#2)/2},#3+5.3) {#4};}
\makeatother

\phchapter{10}{板がたわむ：はりの曲げとたわみ}{本棚の板は、なぜまん中がいちばんたわむ？}{\bklevelrange{4}{5}}{板がたわむ：はりの曲げとたわみ}
\bkmichishirube{5-13,6-3,8-6,力-1,力-5,力-8}{力-11,力-16,波-6}

\begin{goal}
\begin{itemize}[label={\color{bkTheme}●}, leftmargin=1.6\zw, itemsep=0pt, topsep=1pt]
\item \bklevel{3}\enspace 板のたわみが、長さ・厚さ・はば・置く物の重さで、どう変わるかを言える。
\item \bklevel{4}\enspace はりを区切って、せん断力・曲げモーメント・たわみを、表計算の足し算で求められる。
\item \bklevel{5}\enspace つり合いと曲げの式から、$\sigma=\dfrac{My}{I}$ とたわみの微分方程式を導き、代表的なはりのたわみを計算できる。
\end{itemize}
\end{goal}

\phkoushiki{%
\phkx{$R_{\mathrm A}=\dfrac{b}{L}F$、\ $R_{\mathrm B}=\dfrac{a}{L}F$（2点で支えた板が支点から受ける力）}{高校物理}{力-5 で導く}
\phkx{$F_1l_1=F_2l_2$（てこのつり合い）}{小学校理科}{力-5 で導く}
\phkx{$\sum F=0$、\ $\sum M=0$（剛体のつり合い）}{高校物理}{力-5 で導く}
\phkx{$F=kx$（ばねののびは力に比例する）}{中学理科}{力-8 で導く}
}
\noindent はりの曲げとたわみの式は、中学・高校では習いません。でも、その第一歩は、中学・高校で習う「つり合い」の式です。この章では、つり合いの式から出発して、大学の材料力学のはりの式まで登ります。

\begin{hatena}
\par\vspace{3pt}\noindent\begin{minipage}[c]{0.47\linewidth}
本棚の板に本をぎっしり並べると、板が少したわみます。いちばん下がっているのは、板のまん中です。

本はどこにも同じように並んでいるのに、なぜまん中だけが下がるのでしょう。板の厚さを2倍にしたら、たわみは半分になるのかな？
\end{minipage}\hfill
\begin{minipage}[c]{0.50\linewidth}
\centering
\begin{tikzpicture}[x=1mm, y=1mm]
  \fill[bkOrange!25] (0,0) rectangle (4,34);
  \draw[bkOrange!70!black, line width=0.6pt] (0,0) rectangle (4,34);
  \fill[bkOrange!25] (56,0) rectangle (60,34);
  \draw[bkOrange!70!black, line width=0.6pt] (56,0) rectangle (60,34);
  % たわんだ板（大きく見せている）
  \fill[bkOrange!45] (4,17) .. controls (20,12.6) and (40,12.6) .. (56,17) -- (56,15) .. controls (40,10.6) and (20,10.6) .. (4,15) -- cycle;
  \draw[bkOrange!70!black, line width=0.6pt] (4,17) .. controls (20,12.6) and (40,12.6) .. (56,17) -- (56,15) .. controls (40,10.6) and (20,10.6) .. (4,15) -- cycle;
  \draw[bkGray, dashed, line width=0.5pt] (4,17) -- (56,17);
  % 本
  \foreach \x/\c in {6/bkBlue,10/bkRed,14/bkGreen,18/bkPurple,22/bkBlue,26/bkYellow,30/bkRed,34/bkGreen,38/bkBlue,42/bkPurple,46/bkRed,50/bkGreen} {
    \pgfmathsetmacro\yb{17-(4.4*(1-((\x+1.8-30)/26)^2))}
    \fill[\c!55] (\x,\yb) rectangle ++(3.6,14);
    \draw[\c!70!black, line width=0.4pt] (\x,\yb) rectangle ++(3.6,14);}
  \draw[bk arrow, bkInk, line width=0.8pt] (30,6) -- (30,11.2);
  \node[anchor=north, font=\sffamily\footnotesize] at (30,5.6) {まん中がいちばん下がる};
  \node[anchor=south, font=\sffamily\footnotesize, text=bkGray] at (30,34.5) {（たわみを大きくかいた図）};
\end{tikzpicture}
\end{minipage}
\end{hatena}

\noindent 橋・建物の床・クレーンの腕・工場の配管は、どれも「細長い物が、横から力を受けて曲がる」形をしています。このような細長い物を\textbf{はり}といいます。この章では、はりの中で力がどう伝わり、どれだけ曲がるのかを、実験・表計算・式の3つの段で調べます。

\begin{bkfigure}
\begin{tikzpicture}[x=1mm, y=1mm]
  % 橋
  \begin{scope}[shift={(0,0)}]
    \fill[bkGray!40] (2,0) rectangle (6,12) (24,0) rectangle (28,12);
    \draw[bkInk, line width=0.5pt] (2,0) rectangle (6,12) (24,0) rectangle (28,12);
    \fill[bkOrange!40] (0,12) rectangle (30,15); \draw[bkOrange!70!black, line width=0.6pt] (0,12) rectangle (30,15);
    \bkcar{(11,15)}{0.3}{bkBlue}
    \draw[bkBlue!50, line width=0.6pt] (-2,0) -- (32,0);
    \node[anchor=north, font=\sffamily\footnotesize] at (15,-1.5) {橋のけた};
  \end{scope}
  % 床
  \begin{scope}[shift={(36,0)}]
    \fill[bkGray!40] (0,0) rectangle (3,20) (27,0) rectangle (30,20);
    \draw[bkInk, line width=0.5pt] (0,0) rectangle (3,20) (27,0) rectangle (30,20);
    \fill[bkOrange!40] (3,13) rectangle (27,16); \draw[bkOrange!70!black, line width=0.6pt] (3,13) rectangle (27,16);
    \bkperson{(12,16)}{9}{bkRed}
    \bkperson{(19,16)}{8}{bkGreen}
    \node[anchor=north, font=\sffamily\footnotesize] at (15,-1.5) {床のはり};
  \end{scope}
  % クレーン
  \begin{scope}[shift={(72,0)}]
    \fill[bkYellow!60] (2,0) rectangle (6,22); \draw[bkInk, line width=0.5pt] (2,0) rectangle (6,22);
    \fill[bkYellow!60] (2,20) rectangle (30,23); \draw[bkInk, line width=0.5pt] (2,20) rectangle (30,23);
    \draw[bkInk, line width=0.5pt] (27,20) -- (27,11);
    \fill[bkGray!50] (24,6) rectangle (30,11); \draw[bkInk, line width=0.5pt] (24,6) rectangle (30,11);
    \node[anchor=north, font=\sffamily\footnotesize] at (15,-1.5) {クレーンの腕};
  \end{scope}
  % 配管
  \begin{scope}[shift={(108,0)}]
    \draw[bkInk, line width=0.8pt] (-2,22) -- (26,22);
    \foreach \x in {2,22} {\draw[bkInk, line width=0.5pt] (\x,22) -- (\x,13);}
    \fill[bkBlue!30] (-2,10) rectangle (26,13); \draw[bkBlue!70!black, line width=0.6pt] (-2,10) rectangle (26,13);
    \node[anchor=north, font=\sffamily\footnotesize] at (12,-1.5) {つった配管};
  \end{scope}
\end{tikzpicture}
\bkfigcaption{身のまわりのはり。どれも、細長い部材が、横（上下）の向きに力を受けて曲がる。}
\end{bkfigure}

%% ---------------- やってみよう・思い出そう ----------------
\begin{yattemiyou}[定規をたわませてみよう]
\par\vspace{3pt}\noindent\begin{minipage}[c]{0.55\linewidth}
\textbf{用意するもの}\par
30 cm のプラスチックの定規（はば 30 mm・厚さ 3 mm）、同じ高さの本2さつ、500 mL の水の入ったペットボトル、ひも、もう1本の定規
\begin{enumerate}[label={\bkcircled[bkGreen]{\arabic*}}, leftmargin=2.4\zw, itemsep=1pt, topsep=3pt]
\item 本を2さつ立てて、そのはしに定規を寝かせてわたそう。支える所の間の長さ（スパン）を 10 cm にする。
\item 定規のまん中にひもでペットボトルをつるし、まん中が何 mm 下がったかを、もう1本の定規で読もう。
\item スパンを 15 cm、20 cm、25 cm に変えて、同じようにはかろう。
\item ペットボトルをまん中からはしへずらすと、下がり方はどう変わるかな。
\end{enumerate}
\end{minipage}\hfill
\begin{minipage}[c]{0.41\linewidth}
\centering
\begin{tikzpicture}[x=1mm, y=1mm]
  \fill[bkBlue!25] (0,0) rectangle (8,20); \draw[bkBlue!70!black, line width=0.5pt] (0,0) rectangle (8,20);
  \fill[bkBlue!25] (44,0) rectangle (52,20); \draw[bkBlue!70!black, line width=0.5pt] (44,0) rectangle (52,20);
  \draw[bkInk, line width=1.4pt] (5,20.7) .. controls (18,18.6) and (34,18.6) .. (47,20.7);
  \draw[bkInk, line width=0.5pt] (26,19.1) -- (26,13);
  \fill[bkBlue!20] (23,1) rectangle (29,12); \draw[bkBlue!70!black, line width=0.5pt] (23,1) rectangle (29,12);
  \fill[bkBlue!20] (24.4,12) rectangle (27.6,13);
  \draw[bkGray, line width=0.4pt, <->] (8,24) -- (44,24);
  \node[anchor=south, font=\sffamily\footnotesize] at (26,24.4) {スパン};
  \node[anchor=west, font=\sffamily\footnotesize] at (30,8) {水 500 mL};
\end{tikzpicture}
\end{minipage}

\par\vspace{6pt}\noindent 記録の表（ノートに写して書きこもう）
\begin{bkcenter}
\begin{bktable}{colspec={Q[l,m,wd=40mm] *{4}{Q[c,m,wd=15mm]}}, rowsep=2pt}
スパン（cm） & 10 & 15 & 20 & 25 \\
まん中が下がった量（mm） & & & & \\
10 cm のときの何倍か & 1 & & & \\
\end{bktable}
\end{bkcenter}
\end{yattemiyou}

\begin{omoidasou}
\bkrecall{力-5}{回す働き（力のモーメント）＝力×うでの長さ。止まっている物では、力の和も、どの点のまわりのモーメントの和も0。}
\bkrecall{力-8}{応力 $\sigma=\dfrac{F}{A}$、ひずみ $\varepsilon=\dfrac{\Delta L}{L}$。弾性の範囲では $\sigma=E\varepsilon$（$E$ はヤング率）。}
\bkrecall{力-1}{曲がった道に、その場所でぴったり合う円の半径が曲率半径 $\rho$（ロー）。$\rho$ が小さいほど急なカーブ。}
\bkrecall{5-13}{積分すると $+C$ がつく。$C$ は「はじめの状態」の条件で決める。2回積分すると $C$ は2つ。}
\bkrecall{6-3}{$\dfrac{\dd^2y}{\dd x^2}=$（$x$ の式）の形は、積分するだけで解ける。両はしの条件で $C$ を決めることもある。}
\bkrecall{8-6}{細かい面積 $\dd A$ ごとに量をかけて足すのが、面の上の積分（重積分）。}
\end{omoidasou}

%% ================================================================
\bksection{doguu}{物理の道具}

\phstage{me}{たわみを目で見る}{3}

\bksubsection{1}{長さを2倍にすると、たわみは8倍}
\label{ph:力-10:sec1}やってみようの結果を、ミナさんが表にまとめました（ペットボトルはまん中につるした）。

\begin{bkcenter}
\begin{bktable}{colspec={Q[l,m,wd=38mm] *{4}{Q[c,m,wd=17mm]}}, rowsep=1.5pt}
スパン（cm） & 10 & 15 & 20 & 25 \\
まん中が下がった量（mm） & 0.5 & 1.7 & 4.0 & 7.9 \\
10 cm のときの何倍か & 1 & 3.4 & 8.0 & 15.8 \\
\end{bktable}
\bktabcaption{スパンを変えたときの、定規のまん中のたわみ（ミナさんがはかった値）}
\end{bkcenter}
スパンを 10 cm から 20 cm へ2倍にすると、たわみは 0.5 mm から 4.0 mm へ、8倍になりました。$8=2\times2\times2=2^3$ です。15 cm（1.5倍）では $1.5^3=3.375$ で約3.4倍、25 cm（2.5倍）では $2.5^3=15.6$ で約16倍。どれも「長さの3乗」に比例しています。

\noindent\begin{minipage}[c]{0.52\linewidth}
長さが2倍で、なぜ2倍ではなく8倍なのでしょう。理由は3つあります。
\begin{enumerate}[label={\bkcircled{\arabic*}}, leftmargin=2.2\zw, itemsep=1pt, topsep=2pt]
\item うでが2倍になるので、板を曲げる働き（モーメント）が2倍。
\item 曲がる部分の長さが2倍。同じ曲がり方でも、長いほど先が大きく下がる。曲がる角が2倍。
\item その角で進む長さも2倍なので、下がる量はさらに2倍。
\end{enumerate}
$2\times2\times2=8$ 倍です。式でわかるの段（⑧・⑨）で、式からも確かめます。
\end{minipage}\hfill
\begin{minipage}[c]{0.44\linewidth}
\centering
\begin{tikzpicture}[x=1mm, y=1mm]
\begin{axis}[bkplot, width=56mm, height=44mm, xmin=0, xmax=28, ymin=0, ymax=9, xtick={10,15,20,25}, ytick={2,4,6,8}, xlabel={cm}, ylabel={mm}]
  \addplot[bkcurve, domain=0:26, samples=60] {7.88*(x/25)^3};
  \fill[bkRed] (axis cs:10,0.5) circle (1.6pt);
  \fill[bkRed] (axis cs:15,1.7) circle (1.6pt);
  \fill[bkRed] (axis cs:20,4.0) circle (1.6pt);
  \fill[bkRed] (axis cs:25,7.9) circle (1.6pt);
\end{axis}
\end{tikzpicture}
\par\vspace{1pt}{\footnotesize スパンとたわみ（線は「3乗に比例」の曲線）}
\end{minipage}

\pxkhead{置く場所で変わる}
ペットボトルをまん中からはしへずらしていくと、まん中の下がり方は小さくなり、支点の真上ではほとんど0になります。力は同じでも、支点に近い所の力は、板を曲げる「うで」が短いからです。本棚の板でも、重い本をはしに置くと、たわみは小さくなります。

\bksubsection{2}{寝かせるか、立てるか：厚さは3乗でききく}
\label{ph:力-10:sec2}同じ定規を、こんどは立てて（はばの 30 mm をたてにして）わたしてみます。同じペットボトルをつるしても、ほとんど下がりません。計算すると、たわみは寝かせたときの $\dfrac{1}{100}$ です。

\begin{bkfigure}
\begin{tikzpicture}[x=1mm, y=1mm]
  % 寝かせた断面
  \fill[bkBlue!30] (0,10) rectangle (30,13); \draw[bkBlue!70!black, line width=0.6pt] (0,10) rectangle (30,13);
  \draw[bkGray, line width=0.4pt, <->] (0,7) -- (30,7); \node[anchor=north, font=\sffamily\footnotesize] at (15,6.5) {はば 30 mm};
  \draw[bkGray, line width=0.4pt, <->] (33,10) -- (33,13); \node[anchor=west, font=\sffamily\footnotesize] at (34,11.5) {厚さ 3 mm};
  \node[anchor=south, font=\sffamily\small\bfseries] at (15,40) {寝かせる};
  \node[anchor=north, font=\sffamily\footnotesize, align=center] at (15,0) {たわみ 7.9 mm\\（スパン 25 cm）};
  \draw[bk arrow, bkOrange, line width=1.1pt] (15,22) -- (15,14.2);
  \node[anchor=south, font=\sffamily\footnotesize, text=bkOrange!70!black] at (15,22.3) {力の向き};
  % 立てた断面
  \fill[bkRed!25] (70,0) rectangle (73,30); \draw[bkRed!70!black, line width=0.6pt] (70,0) rectangle (73,30);
  \draw[bkGray, line width=0.4pt, <->] (76,0) -- (76,30); \node[anchor=west, font=\sffamily\footnotesize] at (77,15) {高さ};
  \node[anchor=west, font=\sffamily\footnotesize] at (77,10.5) {30 mm};
  \node[anchor=south, font=\sffamily\small\bfseries] at (71.5,40) {立てる};
  \node[anchor=north, font=\sffamily\footnotesize, align=center] at (71.5,-1) {たわみ 0.079 mm\\（100分の1）};
  \draw[bk arrow, bkOrange, line width=1.1pt] (71.5,39.5) -- (71.5,31.2);
  % 説明
  \node[anchor=west, font=\sffamily\small, text width=34mm, align=left] at (97,16) {力の向きに測った\\「高さ」が10倍\\→ たわみは\\$10\times10=100$ 分の1};
\end{tikzpicture}
\bkfigcaption{定規の断面（切り口）。力は上から下へかかる。力の向きの「高さ」が大きいほど、曲がりにくい。}
\end{bkfigure}

\noindent 高さが10倍なのに、なぜ100分の1なのでしょう。もう1つ実験をして考えます。

\pxkhead{2枚重ねるか、くっつけるか}
定規を2枚重ねて（くっつけないで）たわませると、たわみは1枚のときの半分です。2枚がそれぞれ同じ力の半分ずつを受けるからです。ところが、2枚を接着剤でぴったりくっつけて1枚にすると、たわみは1枚のときの $\dfrac{1}{8}$ になります。

\begin{bkfigure}
\begin{tikzpicture}[x=1mm, y=1mm]
  % 1枚
  \begin{scope}[shift={(0,0)}]
    \draw[bkBlue!70!black, line width=0.6pt, fill=bkBlue!25] (0,10) .. controls (12,4) and (24,4) .. (36,10) -- (36,12) .. controls (24,6) and (12,6) .. (0,12) -- cycle;
    \node[anchor=north, font=\sffamily\footnotesize, align=center] at (18,3) {1枚\\たわみ 1};
  \end{scope}
  % 2枚重ね（ずれる）
  \begin{scope}[shift={(44,0)}]
    \draw[bkBlue!70!black, line width=0.6pt, fill=bkBlue!25] (0,12) .. controls (12,8) and (24,8) .. (36,12) -- (36,14) .. controls (24,10) and (12,10) .. (0,14) -- cycle;
    \draw[bkBlue!70!black, line width=0.6pt, fill=bkBlue!25] (0.8,9.6) .. controls (12.8,5.6) and (23.2,5.6) .. (35.2,9.6) -- (35.2,11.6) .. controls (23.2,7.6) and (12.8,7.6) .. (0.8,11.6) -- cycle;
    \draw[bkRed, line width=0.8pt] (0,14) -- (0.8,9.6); \draw[bkRed, line width=0.8pt] (36,14) -- (35.2,9.6);
    \node[anchor=north, font=\sffamily\footnotesize, align=center] at (18,3) {重ねるだけ\\たわみは半分};
    \node[anchor=south, font=\sffamily\footnotesize, text=bkRed!80!black] at (18,15) {はしがずれる};
  \end{scope}
  % くっつける
  \begin{scope}[shift={(88,0)}]
    \draw[bkPurple!70!black, line width=0.6pt, fill=bkPurple!25] (0,12) .. controls (12,10.5) and (24,10.5) .. (36,12) -- (36,16) .. controls (24,14.5) and (12,14.5) .. (0,16) -- cycle;
    \node[anchor=north, font=\sffamily\footnotesize, align=center] at (18,3) {くっつけて1枚\\たわみは8分の1};
    \node[anchor=south, font=\sffamily\footnotesize, text=bkPurple!80!black] at (18,17) {ずれない};
  \end{scope}
\end{tikzpicture}
\bkfigcaption{重ねただけの2枚は、曲げると重なった面ですべって、はしがずれる。くっつけると、ずれることができない。}
\end{bkfigure}

\noindent 重ねただけの2枚を曲げると、上の板の下の面と、下の板の上の面がこすれて、はしがずれます。曲がったはりの外側（下）は長くなろうとし、内側（上）は短くなろうとするからです。くっつけるとずれることができないので、上の板はちぢめられ、下の板はのばされます。材料をのばしたりちぢめたりするには、大きな力がいります。だから、くっつけた板は曲がりにくいのです。

\pxkhead{消しゴムで、はりの中を見る}
\par\vspace{3pt}\noindent\begin{minipage}[c]{0.51\linewidth}
細長い消しゴムの横の面に、ボールペンで等しい間かくのたての線と、まん中に横の線を1本引きます。両はしを持ち上げるようにして、本棚の板と同じ谷の形に曲げると、
\begin{itemize}[leftmargin=1.2\zw, itemsep=1pt, topsep=2pt, label={\color{bkTheme}●}]
\item 上の面では、たての線の間がせまくなる（ちぢむ）。
\item 下の面では、たての線の間が広がる（のびる）。
\item まん中の横の線は、のびもちぢみもしない。
\item たての線は、曲がっても\textbf{まっすぐなまま}で、まん中の横の線と直角のまま。
\end{itemize}
のびもちぢみもしない所を\textbf{中立軸}といいます。\bkidx{ちゅうりつじく}{中立軸}中立軸から遠い所ほど、大きくのびちぢみします。厚い板が曲がりにくいのは、中立軸から遠い所の材料が、大きくのばされたりちぢめられたりして、強くおし返すからです。
\end{minipage}\hfill
\begin{minipage}[c]{0.45\linewidth}
\centering
\begin{tikzpicture}[x=1mm, y=1mm]
  % まっすぐ
  \fill[bkYellow!40] (0,30) rectangle (44,40); \draw[bkInk, line width=0.6pt] (0,30) rectangle (44,40);
  \foreach \x in {4,10,...,40} {\draw[bkBlue, line width=0.5pt] (\x,30) -- (\x,40);}
  \draw[bkRed, line width=0.7pt] (0,35) -- (44,35);
  \node[anchor=south, font=\sffamily\footnotesize] at (22,40.5) {曲げる前};
  % 曲げた（谷の形。中心は上の (22,63)、中立軸の半径 50）
  \begin{scope}[shift={(22,63)}]
    \fill[bkYellow!40] (244:45) arc (244:296:45) -- (296:55) arc (296:244:55) -- cycle;
    \draw[bkInk, line width=0.6pt] (244:45) arc (244:296:45) -- (296:55) arc (296:244:55) -- cycle;
    \foreach \a in {248,254,...,292} {\draw[bkBlue, line width=0.5pt] (\a:45) -- (\a:55);}
    \draw[bkRed, line width=0.7pt] (244:50) arc (244:296:50);
  \end{scope}
  \node[anchor=south, font=\sffamily\footnotesize] at (22,23.3) {ちぢむ};
  \node[anchor=north, font=\sffamily\footnotesize] at (22,7.3) {のびる};
  \node[anchor=west, font=\sffamily\footnotesize, text=bkRed!80!black] at (45.5,16) {中立軸};
  \node[anchor=north, font=\sffamily\footnotesize] at (22,2.5) {両はしを持ち上げて曲げた};
\end{tikzpicture}
\end{minipage}


\begin{phmijika}[段ボールと集成材]
段ボールは、うすい紙2枚の間に波形の紙をはさんで、はりつけたものです。2枚の紙が中立軸から遠くはなれ、しかもずれないので、うすい紙1枚よりずっと曲がりにくくなります。うすい板を何枚もはりあわせた\textbf{集成材}や\textbf{合板}も、はりつけて「ずれない」ようにした材料です。
\end{phmijika}

%% ================================================================
\phstage{te}{表計算でたわみを出す}{3}

\bksubsection{3}{はりの中の力：せん断力と曲げモーメント}
\label{ph:力-10:sec3}本棚の板のどこが、どれだけ強く曲げられているのでしょう。板を、ある場所で（頭の中で）切ってみます。

\noindent\begin{minipage}[c]{0.46\linewidth}
本棚の板（長さ\qty{0.8}{m}）に、本\qty{30}{kg}を一様に並べたとします。本の重さは $30\times9.8=\qty{294}{N}$ で、板\qty{1}{m}あたり $\dfrac{294}{0.8}=\pxkq{367.5}{N/m}$ です。この「1 m あたりの荷重」を\textbf{分布荷重} $w$ といいます。

左右対称なので、両はしの支点は $\dfrac{294}{2}=\qty{147}{N}$ ずつ、板を上へおします（力-5）。これを\textbf{反力} $R$ といいます。
\end{minipage}\hfill
\begin{minipage}[c]{0.50\linewidth}
\centering
\begin{tikzpicture}[x=1mm, y=1mm]
  \fill[bkOrange!35] (0,0) rectangle (46,3); \draw[bkOrange!70!black, line width=0.6pt] (0,0) rectangle (46,3);
  \pxkdist{0}{46}{3}{$w=\pxkq{367.5}{N/m}$}
  \pxkpin{(0,0)} \pxkroller{(46,0)}
  \draw[bk arrow, bkBlue, line width=1.1pt] (0,-16) -- (0,-6.5); \node[anchor=north, font=\sffamily\footnotesize, text=bkBlue!80!black] at (0,-16.3) {$R_{\mathrm A}=\qty{147}{N}$};
  \draw[bk arrow, bkBlue, line width=1.1pt] (46,-16) -- (46,-6.5); \node[anchor=north, font=\sffamily\footnotesize, text=bkBlue!80!black] at (46,-16.3) {$R_{\mathrm B}=\qty{147}{N}$};
  \draw[bkGray, line width=0.4pt, <->] (0,15) -- (46,15); \node[anchor=south, font=\sffamily\footnotesize] at (23,15.3) {$L=\qty{0.8}{m}$};
  \node[anchor=east, font=\sffamily\footnotesize] at (-1.5,1.5) {A}; \node[anchor=west, font=\sffamily\footnotesize] at (47.5,1.5) {B};
\end{tikzpicture}
\end{minipage}

\pxkhead{左の部分のつり合い}
左はしAから $x$ の所で板を切り、左の部分だけを考えます。左の部分は止まっているので、力がつり合っています（力-5）。左の部分にはたらく力は、支点の上向き\qty{147}{N}と、本の下向き $367.5x$（N）です。

\begin{bkfigure}
\begin{tikzpicture}[x=1mm, y=1mm]
  \fill[bkOrange!35] (0,0) rectangle (40,3); \draw[bkOrange!70!black, line width=0.6pt] (0,0) rectangle (40,3);
  \draw[bkRed, line width=1pt, dashed] (40,-3) -- (40,6);
  \pxkdist{0}{40}{3}{本の重さ $367.5x$}
  \pxkpin{(0,0)}
  \draw[bk arrow, bkBlue, line width=1.1pt] (0,-16) -- (0,-6.5); \node[anchor=east, font=\sffamily\footnotesize, text=bkBlue!80!black] at (-1.5,-13) {\qty{147}{N}};
  \draw[bk arrow, bkPurple, line width=1.3pt] (42.5,8) -- (42.5,-5); \node[anchor=west, font=\sffamily\footnotesize, text=bkPurple!80!black] at (44,4) {右の部分がおす力 $Q$};
  \draw[bkRed, line width=1.2pt, -{Stealth[length=5pt]}] (46,-9) arc (-60:60:6);
  \node[anchor=west, font=\sffamily\footnotesize, text=bkRed!80!black] at (51,-6) {右の部分が回そうとする働き $M$};
  \draw[bkGray, line width=0.4pt, <->] (0,-21) -- (40,-21); \node[anchor=north, font=\sffamily\footnotesize] at (20,-21.4) {$x$};
\end{tikzpicture}
\bkfigcaption{切った所から左の部分。切り口には、右の部分から、上下の力 $Q$ と、回す働き $M$ がはたらく。}
\end{bkfigure}

\noindent 支点の力と本の重さは、ふつう等しくありません。それでも左の部分が落ちないのは、切り口で右の部分が上下の力をおよぼしているからです。この力を\textbf{せん断力} $Q$ といいます。\bkidx{せんだんりょく}{せん断力}この本では、切り口より左にある上向きの力の和を、そのまま $Q$ とよびます。
\[ Q=147-367.5x \]
また、左の部分の力は、切り口のまわりに左の部分を回そうとします。それを止めているのが、右の部分が切り口にあたえる回す働きです。これを\textbf{曲げモーメント} $M$ といいます。\bkidx{まげもーめんと}{曲げモーメント}左の部分の力が、切り口のまわりに回す働き（力×うで）を合計したものが $M$ です。

\noindent\begin{minipage}[c]{0.56\linewidth}
\begin{itemize}[leftmargin=1.2\zw, itemsep=2pt, topsep=0pt, label={\color{bkTheme}●}]
\item 支点の\qty{147}{N}：うでは $x$。回す働きは $147x$。
\item 本の重さ $367.5x$：本は $0$ から $x$ まで一様に並ぶので、重さのまん中 $\dfrac{x}{2}$ に集まっていると考えてよい（力-5 の重心）。うでは $\dfrac{x}{2}$。回す働きは逆向きで $367.5x\times\dfrac{x}{2}$。
\end{itemize}
\end{minipage}\hfill
\begin{minipage}[c]{0.40\linewidth}
\[ M=147x-\frac{367.5}{2}x^2 \]
{\small 板を谷の形（下に凸）に曲げる向きを、$M$ のプラスと決めます。}
\end{minipage}

\bksubsection{4}{表計算：足し算でせん断力図・曲げモーメント図をかく}
\label{ph:力-10:sec4}式を使わなくても、足し算だけで $Q$ と $M$ が求められます。板を\qty{0.1}{m}ずつ8つの区間に分け、左から順に足していきます。
\begin{enumerate}[label={\bkcircled{\arabic*}}, leftmargin=2.2\zw, itemsep=1pt, topsep=2pt]
\item $Q$ は、左はしで反力の\qty{147}{N}。区間を1つ進むたびに、その区間の本の重さ $367.5\times0.1=\qty{36.75}{N}$ だけへる。
\item $M$ は、左はしで0。区間を1つ進むたびに、（その区間のまん中の $Q$）×\qty{0.1}{m} だけふえる。
\end{enumerate}
②がなぜ正しいかは、⑥で式から確かめます。いまは「$Q$ が正の所では $M$ がふえ、$Q$ が負の所では $M$ がへる」と読んでおきます。

\begin{bkcenter}
\bksheet{4}{24}{%
1 & $x$（m） & $Q$（N） & 区間のまん中の $Q$（N） & $M$（N·m） \\
2 & 0.0 & 147.00 & & 0.00 \\
3 & 0.1 & 110.25 & 128.625 & 12.86 \\
4 & 0.2 & 73.50 & 91.875 & 22.05 \\
5 & 0.3 & 36.75 & 55.125 & 27.56 \\
6 & 0.4 & 0.00 & 18.375 & 29.40 \\
7 & 0.5 & $-36.75$ & $-18.375$ & 27.56 \\
8 & 0.6 & $-73.50$ & $-55.125$ & 22.05 \\
9 & 0.7 & $-110.25$ & $-91.875$ & 12.86 \\
10 & 0.8 & $-147.00$ & $-128.625$ & 0.00 \\}
\bktabcaption{B3 は \texttt{=B2-367.5*0.1}、C3 は \texttt{=(B2+B3)/2}、D3 は \texttt{=D2+C3*0.1}。下の行へコピーする。}
\end{bkcenter}

\begin{bkfigure}
\begin{tikzpicture}[x=1mm, y=1mm]
\begin{axis}[bkplot, axis lines=left, name=Q, width=56mm, height=44mm, xmin=0, xmax=0.88, ymin=-170, ymax=170, xtick={0.2,0.4,0.6,0.8}, ytick={-147,0,147}, xlabel={$x$（m）}, ylabel={せん断力 $Q$（N）}]
  \draw[bkInk, line width=0.6pt] (axis cs:0,0) -- (axis cs:0.88,0);
  \fill[bkBlue!15] (axis cs:0,0) -- (axis cs:0,147) -- (axis cs:0.8,-147) -- (axis cs:0.8,0) -- cycle;
  \draw[bkBlue, line width=1.6pt] (axis cs:0,0) -- (axis cs:0,147) -- (axis cs:0.8,-147) -- (axis cs:0.8,0);
\end{axis}
\begin{axis}[bkplot, axis lines=left, at={(68mm,0)}, width=56mm, height=44mm, xmin=0, xmax=0.88, ymin=0, ymax=36, xtick={0.2,0.4,0.6,0.8}, ytick={10,20,29.4}, xlabel={$x$（m）}, ylabel={曲げモーメント $M$（N·m）}]
  \addplot[bkRed, line width=1.6pt, fill=bkRed!12, domain=0:0.8, samples=60] {147*x-183.75*x^2} \closedcycle;
\end{axis}
\end{tikzpicture}
\bkfigcaption{本棚の板のせん断力図（左）と曲げモーメント図（右）。$Q$ が0になる所（まん中）で、$M$ がいちばん大きい（\pxkq{29.4}{N.m}）。}
\end{bkfigure}

\noindent $Q$ の図は、左はしの\qty{147}{N}から右はしの $-\qty{147}{N}$ まで、まっすぐにへる線です。$M$ の図は、まん中で山になる放物線です。このように、はりの場所ごとの $Q$ と $M$ をグラフにしたものを、\textbf{せん断力図}・\textbf{曲げモーメント図}といいます。設計の図面の横には、よくこの2つの図がかいてあります。

\bksubsection{5}{表計算：曲がり方を2回足して、たわみを出す}
\label{ph:力-10:sec5}板をいちばん強く曲げているのは、$M$ がいちばん大きいまん中です。では、板はどれだけ下がるのでしょう。

①・②の実験から、板の曲がり方は、曲げる働き $M$ に比例し、材料のかたさ（ヤング率 $E$）と断面の形で決まる「曲げにくさ」に反比例することがわかっています。この「曲げにくさ」を $EI$（\textbf{曲げ剛性}）と書きます。\bkidx{まげごうせい}{曲げ剛性}$I$ は断面の形で決まる量で、はば $b$、高さ $h$ の長方形では $I=\dfrac{bh^3}{12}$ です（⑦で導く。高さの3乗が、②の100倍の理由）。

\[ \text{曲がり方}=\frac{M}{EI}\qquad\Bigl(\text{単位は }\frac{1}{\mathrm m}\text{。曲率半径 }\rho\text{ の逆数}\Bigr) \]

\noindent 本棚の板（すぎ、$E=\qty{7}{GPa}$、はば\qty{0.25}{m}、厚さ\qty{0.018}{m}）では、
\[ I=\frac{0.25\times0.018^3}{12}=\qty{1.215e-7}{m^4} \]
\[ EI=7\times10^9\times1.215\times10^{-7}=\pxkq{850.5}{N.m^2} \]

\pxkhead{まん中から右へ、2回足す}
板は左右対称なので、まん中では板は水平です。まん中から右へ\qty{0.05}{m}ずつ進みながら、
\begin{enumerate}[label={\bkcircled{\arabic*}}, leftmargin=2.2\zw, itemsep=1pt, topsep=2pt]
\item 曲がり方 $\dfrac{M}{EI}$ を足していくと、板の\textbf{傾き}（1 m 進むと何 mm 上がるか）が出る。
\item 傾きを足していくと、まん中からはかって板が\textbf{上がった高さ}が出る。
\end{enumerate}
右はし（支点）まで来たときの「上がった高さ」が、まん中が支点より下がっている量、つまり\textbf{たわみ}です。足すときは、前の行とこの行の平均に\qty{0.05}{m}をかけます（台形で面積を求めるのと同じ。数学 5-10）。

\begin{bkcenter}
\bksheet{5}{21}{%
1 & まん中から（m） & $M$（N·m） & 曲がり方（1/km） & 傾き（mm/m） & 上がった高さ（mm） \\
2 & 0.00 & 29.400 & 34.57 & 0.000 & 0.0000 \\
3 & 0.05 & 28.941 & 34.03 & 1.715 & 0.0429 \\
4 & 0.10 & 27.562 & 32.41 & 3.376 & 0.1701 \\
5 & 0.15 & 25.266 & 29.71 & 4.929 & 0.3777 \\
6 & 0.20 & 22.050 & 25.93 & 6.319 & 0.6590 \\
7 & 0.25 & 17.916 & 21.07 & 7.494 & 1.0043 \\
8 & 0.30 & 12.862 & 15.12 & 8.399 & 1.4016 \\
9 & 0.35 & 6.891 & 8.10 & 8.980 & 1.8361 \\
10 & 0.40 & 0.000 & 0.00 & 9.182 & 2.2901 \\}
\bktabcaption{C2 は \texttt{=B2/850.5*1000}、D3 は \texttt{=D2+(C2+C3)/2*0.05}、E3 は \texttt{=E2+(D2+D3)/2*0.05}（1/km は 1000 m あたり。mm/m×m＝mm）}
\end{bkcenter}
右はしで約\qty{2.3}{mm}。これが本棚の板のたわみです。区間をもっと細かくすると、答えは\qty{2.305}{mm}に近づきます。プログラムで、区間を1000にしてみましょう。

\begin{phcode}
L, w, EI = 0.8, 367.5, 850.5
def M(x):                     \# 左はしから x の所の曲げモーメント
    return w*L/2*x - w*x*x/2
n = 1000
ds = (L/2)/n
k0, th, up = M(L/2)/EI, 0.0, 0.0
for i in range(1, n+1):
    k1 = M(L/2 + i*ds)/EI     \# 曲がり方
    th1 = th + (k0 + k1)/2*ds   \# 傾き
    up = up + (th + th1)/2*ds   \# 上がった高さ
    th, k0 = th1, k1
print(round(up*1000, 3), "mm")
\end{phcode}
\begin{phcodeout}
2.305 mm
\end{phcodeout}

\noindent この足し算で、\bkhatenaref の答えの半分が見えてきました。曲がり方は $M$ に比例し、$M$ はまん中で最大です。まん中で水平な板が、支点に向かって少しずつ曲がり上がっていくので、両はしの支点にのせると、\textbf{まん中がいちばん低く}なります。

\phkokomade

%% ================================================================
\phstage{shiki}{つり合いと曲げの式から導く}{5}

\bksubsection{6}{小さな部分のつり合い：$\dfrac{\dd Q}{\dd x}=-w$、\ $\dfrac{\dd M}{\dd x}=Q$}
\label{ph:力-10:sec6}④の足し算のきまりを、式で確かめます。はりから、長さ $\dd x$ の小さな部分を切り出します。

\noindent\begin{minipage}[c]{0.57\linewidth}
\centering
\begin{tikzpicture}[x=1mm, y=1mm]
  \fill[bkOrange!30] (10,0) rectangle (40,10); \draw[bkOrange!70!black, line width=0.7pt] (10,0) rectangle (40,10);
  \pxkdist{10}{40}{10}{$w\,\dd x$}
  \draw[bk arrow, bkBlue, line width=1.2pt] (6,-6) -- (6,12); \node[anchor=east, font=\sffamily\footnotesize, text=bkBlue!80!black] at (5,3) {$Q$};
  \draw[bk arrow, bkBlue, line width=1.2pt] (44,16) -- (44,-2); \node[anchor=west, font=\sffamily\footnotesize, text=bkBlue!80!black] at (45,10) {$Q+\dd Q$};
  \draw[bkRed, line width=1.1pt, -{Stealth[length=4.5pt]}] (5.2,6.9) ++(250:6.5) arc (250:110:6.5); \node[anchor=east, font=\sffamily\footnotesize, text=bkRed!80!black] at (-1,-4) {$M$};
  \draw[bkRed, line width=1.1pt, -{Stealth[length=4.5pt]}] (47,-3) arc (-70:70:6.5); \node[anchor=west, font=\sffamily\footnotesize, text=bkRed!80!black] at (51,-4) {$M+\dd M$};
  \draw[bkGray, line width=0.4pt, <->] (10,-8) -- (40,-8); \node[anchor=north, font=\sffamily\footnotesize] at (25,-8.4) {$\dd x$};
\end{tikzpicture}
\end{minipage}\hfill
\begin{minipage}[c]{0.40\linewidth}
左の切り口には、左の部分から $Q$（上向き）と $M$、右の切り口には、右の部分から $Q+\dd Q$（下向き）と $M+\dd M$ がはたらきます（左と右で、作用と反作用なので向きが逆）。上には荷重 $w\,\dd x$ がのっています。
\end{minipage}

\pxkhead{上下の力のつり合い}
\[ Q-(Q+\dd Q)-w\,\dd x=0\qquad\text{より}\qquad \frac{\dd Q}{\dd x}=-w \]
\pxkhead{右の切り口のまわりのモーメントのつり合い}
左の切り口の $Q$ のうでは $\dd x$、荷重 $w\,\dd x$ のうでは $\dfrac{\dd x}{2}$ です。
\[ M+Q\,\dd x-w\,\dd x\cdot\frac{\dd x}{2}-(M+\dd M)=0 \]
$\dd x$ はとても小さいので、$\dd x$ の2乗の項は、ほかの項よりずっと小さくなって消えます（数学 5-6 の「3つめの項」と同じ）。
\pheq{\frac{\dd M}{\dd x}=Q}{N ＝ N·m ÷ m}

\begin{phyomi}[$\dfrac{\dd Q}{\dd x}=-w$、\ $\dfrac{\dd M}{\dd x}=Q$]
$Q$ のグラフの傾きは「1 m あたりの荷重」に $-$ をつけたもの。$M$ のグラフの傾きは $Q$。だから、$Q=0$ の所で $M$ のグラフは水平になり、$M$ は最大か最小になる（数学 5-7）。④の表で、$Q$ が0のまん中で $M$ が最大だったのは、このためです。
\end{phyomi}
\begin{phmijika}
荷物を1つだけのせた板では、荷物のない所は $w=0$ なので、$Q$ は一定（$Q$ 図は水平）、$M$ はまっすぐな線（$M$ 図は三角形）になります。
\end{phmijika}

\pxkhead{1つの荷物をのせた板：支点の力を導く}
\label{ph:力-10:hanryoku}長さ $L$ の板の、左はしAから $a$、右はしBから $b$（$a+b=L$）の所に、荷物 $F$ をのせます。荷物のない所は $w=0$ なので、$Q$ は一定です。A〜荷物の間の $Q$ は $R_{\mathrm A}$ のまま、荷物をこえると $F$ だけへって $R_{\mathrm A}-F$ になります。$M$ は $Q$ を足したもので、はしでは0です（はしの支点は、板を回す働きをおよぼさない）。
\[ M(L)=R_{\mathrm A}a+(R_{\mathrm A}-F)b=R_{\mathrm A}L-Fb=0 \]
\[ R_{\mathrm A}=\frac{b}{L}F,\qquad R_{\mathrm B}=F-R_{\mathrm A}=\frac{a}{L}F \]
これは、Bのまわりのモーメントのつり合い $R_{\mathrm A}L=Fb$、つまり\textbf{てこのつり合い}（力-5）そのものです。荷物に近い支点ほど、大きな力で板を支えます。

\Needspace{14\baselineskip}
\begin{phflow}
\phflowrow{$\dfrac{\dd M}{\dd x}=Q$、\ $\dfrac{\dd Q}{\dd x}=-w$}{荷重は $F$ だけ\\両はしで $M=0$}{$R_{\mathrm A}L=Fb$\\（てこのつり合い）}
\phflowrow{$R_{\mathrm A}L=Fb$}{荷物がまん中\\（$a=b$）}{$R_{\mathrm A}=R_{\mathrm B}=\dfrac{F}{2}$}
\end{phflow}

\bksubsection{7}{曲げ応力 $\sigma=\dfrac{My}{I}$ と断面二次モーメント}
\label{ph:力-10:sec7}②の消しゴムの実験から、「曲げても断面は平らなまま」とします（\textbf{平面保持}の仮定）。このとき、はりの中のひずみと応力が決まります。

\begin{bkfigure}
\begin{tikzpicture}[x=1mm, y=1mm]
  % 曲がった小さな部分（中心 O は上、中立軸の半径 36）
  \coordinate (O) at (30,50);
  \fill[bkYellow!35] ($(O)+(245:30)$) arc (245:295:30) -- ($(O)+(295:42)$) arc (295:245:42) -- cycle;
  \draw[bkInk, line width=0.7pt] ($(O)+(245:30)$) arc (245:295:30) -- ($(O)+(295:42)$) arc (295:245:42) -- cycle;
  \draw[bkRed, line width=0.9pt] ($(O)+(245:36)$) arc (245:295:36);
  \draw[bkBlue, line width=0.9pt] ($(O)+(245:40)$) arc (245:295:40);
  \draw[bkGray, dashed, line width=0.4pt] (O) -- ($(O)+(245:44)$); \draw[bkGray, dashed, line width=0.4pt] (O) -- ($(O)+(295:44)$);
  \fill[bkInk] (O) circle (0.8); \node[anchor=south, font=\sffamily\footnotesize] at ($(O)+(0,0.8)$) {曲がりの中心};
  \draw[bkInk, line width=0.5pt] ($(O)+(245:8)$) arc (245:295:8); \node[font=\sffamily\footnotesize] at ($(O)+(270:11)$) {$\dd\theta$};
  \draw[bkGray, line width=0.4pt, <->] ($(O)+(238:0.5)$) -- ($(O)+(238:36)$); \node[anchor=east, font=\sffamily\footnotesize] at ($(O)+(236:22)$) {$\rho$};
  \node[anchor=west, font=\sffamily\footnotesize, text=bkRed!80!black] at ($(O)+(296:36)+(1,0)$) {中立軸：長さ $\rho\,\dd\theta$};
  \node[anchor=west, font=\sffamily\footnotesize, text=bkBlue!80!black] at ($(O)+(296:40)+(1,-1.5)$) {$y$ だけ下：$(\rho+y)\,\dd\theta$};
  % 応力の分布
  \begin{scope}[shift={(108,16)}]
    \draw[bkInk, line width=0.7pt] (0,-8) -- (0,8);
    \draw[bkGray, line width=0.4pt, dashed] (-14,0) -- (14,0); \node[anchor=west, font=\sffamily\footnotesize, text=bkRed!80!black] at (14.5,0) {中立軸};
    \fill[bkRed!20] (0,0) -- (12,-8) -- (0,-8) -- cycle; \fill[bkBlue!20] (0,0) -- (-12,8) -- (0,8) -- cycle;
    \foreach \yy in {-7,-5,-3,-1} {\draw[-{Stealth[length=3pt]}, bkRed, line width=0.7pt] (0,\yy) -- ({-1.5*\yy},\yy);}
    \foreach \yy in {1,3,5,7} {\draw[-{Stealth[length=3pt]}, bkBlue, line width=0.7pt] (0,\yy) -- ({-1.5*\yy},\yy);}
    \node[anchor=south, font=\sffamily\footnotesize, text=bkBlue!80!black] at (-6,9) {おしちぢめ};
    \node[anchor=north, font=\sffamily\footnotesize, text=bkRed!80!black] at (6,-9) {引っぱり};
    \node[anchor=north, font=\sffamily\footnotesize] at (0,-14) {断面にかかる応力};
  \end{scope}
\end{tikzpicture}
\bkfigcaption{谷の形に曲がったはりの、長さ $\dd x$ の部分。中立軸より $y$ だけ下の線は、中立軸より長い。応力は $y$ に比例する。}
\end{bkfigure}

\pxkhead{（ア）ひずみ：中立軸からの距離に比例}
曲がりの中心からはかった中立軸の半径を $\rho$、この部分が見こむ角を $\dd\theta$ とします。中立軸の長さは $\rho\,\dd\theta$（数学 4-2 の $s=r\theta$）で、のびちぢみしません。中立軸から $y$ だけ下の線の長さは $(\rho+y)\,\dd\theta$ です。のびは $y\,\dd\theta$ なので、
\[ \varepsilon=\frac{y\,\dd\theta}{\rho\,\dd\theta}=\frac{y}{\rho} \]
\pxkhead{（イ）応力：$\sigma=E\varepsilon$（力-8）}
\[ \sigma=E\frac{y}{\rho} \]
中立軸の下（$y>0$）は引っぱられ、上（$y<0$）はおしちぢめられます。
\pxkhead{（ウ）中立軸は、断面の重心を通る}
はりを横から曲げるだけで、軸の向きには引っぱっていないので、断面の応力をすべて足すと0です。
\[ \int\sigma\,\dd A=\frac{E}{\rho}\int y\,\dd A=0\qquad\text{より}\qquad \int y\,\dd A=0 \]
$\displaystyle\int y\,\dd A=0$ は「断面の重心が $y=0$ にある」ということです（力-5 の重心の式と同じ形）。だから、長方形でも円でも、中立軸は断面のまん中を通ります。
\pxkhead{（エ）応力が回す働きの合計が $M$}
小さな面積 $\dd A$ の応力の力 $\sigma\,\dd A$ は、中立軸のまわりに、うで $y$ で回す働き $\sigma y\,\dd A$ をもちます。これを全部足したものが、曲げモーメント $M$ です。
\[ M=\int\sigma y\,\dd A=\frac{E}{\rho}\int y^2\,\dd A \]
ここで出てきた $\displaystyle\int y^2\,\dd A$ を\textbf{断面二次モーメント} $I$ といいます。\bkidx{だんめんにじもーめんと}{断面二次モーメント}
\pheq{I=\int y^2\,\dd A,\qquad \frac{1}{\rho}=\frac{M}{EI}}{1/m ＝ N·m ÷（Pa × m$^4$）}
⑤の「曲がり方 $\dfrac{M}{EI}$」は、この式から出てきたものです。$\dfrac{E}{\rho}=\dfrac{M}{I}$ を（イ）に入れると、
\pheq{\sigma=\frac{My}{I}}{Pa ＝ N·m × m ÷ m$^4$}
\label{ph:力-10:sigma}
\begin{phyomi}[$\sigma=\dfrac{My}{I}$]
曲げ応力は、曲げモーメント $M$ に比例し、中立軸からの距離 $y$ に比例する。断面のいちばん外側（上と下の面）で最大になる。$I$ が大きい断面ほど、同じ $M$ でも応力が小さい。
\end{phyomi}

\pxkhead{長方形の $I=\dfrac{bh^3}{12}$}
はば $b$、高さ $h$ の長方形では、中立軸からの高さ $y$ の所に、はば $b$、厚さ $\dd y$ の細い帯をとると $\dd A=b\,\dd y$ です。
\[ I=\int_{-h/2}^{h/2}y^2\,b\,\dd y=b\Bigl[\frac{y^3}{3}\Bigr]_{-h/2}^{h/2}=b\Bigl(\frac{h^3}{24}+\frac{h^3}{24}\Bigr)=\frac{bh^3}{12} \]
\begin{phbaisuru}{高さ $h$}
$I$ は $2^3=8$ 倍になり、同じ $M$ での曲がり方は $\dfrac{1}{8}$ になる。②で2枚をくっつけた板は、高さが2倍の1枚の板なので、たわみが $\dfrac{1}{8}$ でした。定規を立てると高さは10倍、はばは $\dfrac{1}{10}$ なので、$I$ は $10^3\times\dfrac{1}{10}=100$ 倍です。
\end{phbaisuru}
\begin{phbaisuru}{はば $b$}
$I$ は2倍、曲がり方は半分。重ねただけの2枚は、はばが2倍の板と同じです（たわみは $\dfrac{1}{2}$）。
\end{phbaisuru}

\pxkhead{いろいろな断面の $I$}
円の $I$ は、円の中心からの距離 $r$ を使って求めます。$x^2+y^2=r^2$ なので、$\displaystyle\int x^2\,\dd A+\int y^2\,\dd A=\int r^2\,\dd A$。円は、どの向きから見ても同じ形なので、左の2つの積分は等しく、どちらも $\displaystyle\frac{1}{2}\int r^2\,\dd A$ です。半径 $r$ の所の、はば $\dd r$ の細い輪の面積は $2\pi r\,\dd r$ なので（数学 8-6）、
\[ \int r^2\,\dd A=\int_0^{d/2}r^2\cdot2\pi r\,\dd r=2\pi\Bigl[\frac{r^4}{4}\Bigr]_0^{d/2}=\frac{\pi d^4}{32},\qquad I=\frac{\pi d^4}{64} \]
（$\displaystyle\int r^2\,\dd A$ は、力-11 のねじりで、もう一度使います。）

\begin{bkcenter}
\begin{bktable}{colspec={Q[c,m,wd=25mm] Q[c,m,wd=25mm] Q[c,m,wd=34mm] Q[c,m,wd=30mm]}, rowsep=2pt}
\begin{tikzpicture}[x=1mm, y=1mm]\fill[bkBlue!25] (0,0) rectangle (10,16); \draw[bkBlue!70!black, line width=0.6pt] (0,0) rectangle (10,16); \draw[bkRed, dashed, line width=0.5pt] (-2,8) -- (12,8); \node[anchor=north, font=\sffamily\footnotesize] at (5,-0.5) {$b$}; \node[anchor=west, font=\sffamily\footnotesize] at (10.5,13) {$h$};\end{tikzpicture}
& \begin{tikzpicture}[x=1mm, y=1mm]\fill[bkBlue!25] (0,0) circle (8); \draw[bkBlue!70!black, line width=0.6pt] (0,0) circle (8); \draw[bkRed, dashed, line width=0.5pt] (-10,0) -- (10,0); \node[anchor=north, font=\sffamily\footnotesize] at (0,-8.5) {直径 $d$};\end{tikzpicture}
& \begin{tikzpicture}[x=1mm, y=1mm]\fill[bkBlue!25, even odd rule] (0,0) circle (8) (0,0) circle (6); \draw[bkBlue!70!black, line width=0.6pt] (0,0) circle (8) (0,0) circle (6); \draw[bkRed, dashed, line width=0.5pt] (-10,0) -- (10,0); \node[anchor=north, font=\sffamily\footnotesize] at (0,-8.5) {外径 $D$・内径 $d$};\end{tikzpicture}
& \begin{tikzpicture}[x=1mm, y=1mm]\fill[bkBlue!25] (0,0) rectangle (12,2) (5.3,2) rectangle (6.7,14) (0,14) rectangle (12,16); \draw[bkBlue!70!black, line width=0.6pt] (0,0) -- (12,0) -- (12,2) -- (6.7,2) -- (6.7,14) -- (12,14) -- (12,16) -- (0,16) -- (0,14) -- (5.3,14) -- (5.3,2) -- (0,2) -- cycle; \draw[bkRed, dashed, line width=0.5pt] (-2,8) -- (14,8);\end{tikzpicture} \\
長方形 & 円 & パイプ & H形（例題8） \\
$I=\dfrac{bh^3}{12}$ & $I=\dfrac{\pi d^4}{64}$ & $I=\dfrac{\pi(D^4-d^4)}{64}$ & {外の長方形から\\欠けた部分をひく} \\
$Z=\dfrac{bh^2}{6}$ & $Z=\dfrac{\pi d^3}{32}$ & $Z=\dfrac{\pi(D^4-d^4)}{32D}$ & $Z=\dfrac{I}{h/2}$ \\
\end{bktable}
\bktabcaption{よく使う断面の断面二次モーメント $I$ と断面係数 $Z$（赤い点線が中立軸）}
\end{bkcenter}

\pxkhead{いちばん大きい応力と断面係数}
応力がいちばん大きいのは、中立軸からいちばん遠い面（$y=y_{\max}$）です。$Z=\dfrac{I}{y_{\max}}$ とおくと、
\pheq{\sigma_{\max}=\frac{M}{Z}}{Pa ＝ N·m ÷ m$^3$}
$Z$ を\textbf{断面係数}といいます。長方形なら $y_{\max}=\dfrac{h}{2}$ なので $Z=\dfrac{bh^3}{12}\div\dfrac{h}{2}=\dfrac{bh^2}{6}$ です。たわみ（かたさ）は $I$ で、こわれにくさ（強さ）は $Z$ で決まります。

\bksubsection{8}{たわみの微分方程式 $EI\dfrac{\dd^2y}{\dd x^2}=-M$}
\label{ph:力-10:sec8}はりの軸が、場所 $x$ でどれだけ下がったかを $y(x)$ とします（下向きを正。⑦の断面の中の $y$ と同じ下向きです）。力-1 ⑤で、曲がった線の曲がり方 $\dfrac{1}{\rho}$ は、傾きが小さいときは $\dfrac{\dd^2y}{\dd x^2}$ の大きさで表せました。

\noindent\begin{minipage}[c]{0.47\linewidth}
本棚の板のように谷の形（下に凸）に曲がるとき、$M>0$ です。$y$ を下向きにはかると、谷の形の線では、$y$ は「まん中で最大」の山の形のグラフになり、$\dfrac{\dd^2y}{\dd x^2}<0$ です（数学 5-7）。そこで $-$ をつけて、
\end{minipage}\hfill
\begin{minipage}[c]{0.49\linewidth}
\centering
\begin{tikzpicture}[x=1mm, y=1mm]
  \draw[bkGray, line width=0.6pt, -{Stealth[length=4pt]}] (0,0) -- (52,0) node[anchor=west, font=\sffamily\footnotesize] {$x$};
  \draw[bkGray, line width=0.6pt, -{Stealth[length=4pt]}] (0,0) -- (0,-16) node[anchor=north, font=\sffamily\footnotesize] {$y$（下が正）};
  \draw[bkBlue, line width=1.5pt] (0,0) .. controls (15,-11) and (31,-11) .. (46,0);
  \draw[bkGray, line width=0.4pt, <->] (23,0) -- (23,-8.2); \node[anchor=west, font=\sffamily\footnotesize] at (24,-4) {$\delta$};
  \pxkpin{(0,0)} \pxkroller{(46,0)}
\end{tikzpicture}
\end{minipage}
\pheq{EI\frac{\dd^2y}{\dd x^2}=-M}{N·m ＝ N·m$^2$ × 1/m}
\label{ph:力-10:taw}
と書きます。これが\textbf{たわみの微分方程式}です。⑥の $\dfrac{\dd M}{\dd x}=Q$、$\dfrac{\dd Q}{\dd x}=-w$ と合わせると、$EI\dfrac{\dd^4y}{\dd x^4}=w$ とも書けます。

\pxkhead{解き方：2回積分して、両はしの条件で $C$ を決める}
本棚の板（単純支持・等分布荷重）では、③で $M=\dfrac{w}{2}(Lx-x^2)$ でした。
\begin{enumerate}[label={\bkcircled{\arabic*}}, leftmargin=2.2\zw, itemsep=3pt, topsep=2pt]
\item 1回積分（数学 5-11）：$EI\dfrac{\dd y}{\dd x}=-\dfrac{w}{2}\Bigl(\dfrac{Lx^2}{2}-\dfrac{x^3}{3}\Bigr)+C_1$
\item もう1回積分：$EIy=-\dfrac{w}{2}\Bigl(\dfrac{Lx^3}{6}-\dfrac{x^4}{12}\Bigr)+C_1x+C_2$
\item 左はしで下がらない（$x=0$ で $y=0$）：$C_2=0$
\item 右はしでも下がらない（$x=L$ で $y=0$）：$-\dfrac{w}{2}\Bigl(\dfrac{L^4}{6}-\dfrac{L^4}{12}\Bigr)+C_1L=0$ より $C_1=\dfrac{wL^3}{24}$
\end{enumerate}
\[ y=\frac{w}{24EI}\bigl(L^3x-2Lx^3+x^4\bigr) \]
数学 5-13 では、$C$ を「はじめの時刻の状態」で決めました。はりでは、$C$ を「\textbf{両はしの状態}」で決めます。これを\textbf{境界条件}といいます（数学 6-3）。

\pxkhead{\bkhatenaref の答え}
$\dfrac{\dd y}{\dd x}=0$ になるのは、左右対称なのでまん中 $x=\dfrac{L}{2}$ です。そこで $y$ が最大になります。
\[ \delta=y\Bigl(\frac{L}{2}\Bigr)=\frac{w}{24EI}\Bigl(\frac{L^4}{2}-\frac{L^4}{4}+\frac{L^4}{16}\Bigr)=\frac{5wL^4}{384EI} \]
本棚の板では $\delta=\dfrac{5\times367.5\times0.8^4}{384\times850.5}=\qty{2.30e-3}{m}$、約\qty{2.3}{mm}。⑤の表計算と同じです。まん中がいちばんたわむのは、①曲げモーメントがまん中で最大で、板がいちばん強く曲がっていること、②両はしが支点で止められ、左右対称なのでまん中で傾きが0になり、そこが谷の底になること、の2つのためです。

\begin{bkfigure}
\begin{tikzpicture}[x=1mm, y=1mm]
\begin{axis}[bkplot, width=120mm, height=40mm, xmin=0, xmax=0.86, ymin=-2.6, ymax=0.4, y dir=normal, xtick={0.2,0.4,0.6,0.8}, ytick={-2.3,-1}, yticklabels={$2.3$,$1$}, xlabel={$x$（m）}, ylabel={下がった量（mm）}]
  \addplot[bkBlue, line width=1.7pt, domain=0:0.8, samples=80] {-367.5/(24*850.5)*(0.512*x-1.6*x^3+x^4)*1000};
  \fill[bkRed] (axis cs:0,0) circle (1.5pt);
  \fill[bkRed] (axis cs:0.1,-0.87) circle (1.5pt);
  \fill[bkRed] (axis cs:0.2,-1.599) circle (1.5pt);
  \fill[bkRed] (axis cs:0.3,-2.079) circle (1.5pt);
  \fill[bkRed] (axis cs:0.4,-2.247) circle (1.5pt);
  \fill[bkRed] (axis cs:0.5,-2.079) circle (1.5pt);
  \fill[bkRed] (axis cs:0.6,-1.599) circle (1.5pt);
  \fill[bkRed] (axis cs:0.7,-0.87) circle (1.5pt);
  \fill[bkRed] (axis cs:0.8,0) circle (1.5pt);
  \node[anchor=west, font=\sffamily\footnotesize, text=bkRed!80!black] at (axis cs:0.52,-0.6) {●は表計算（0.1 m 区切り）};
\end{axis}
\end{tikzpicture}
\bkfigcaption{本棚の板のたわみの形（線は式、●は表計算）。グラフは、下がる向きを下にかいた。}
\end{bkfigure}

\begin{phyomi}[$\delta=\dfrac{5wL^4}{384EI}$]
分子は「荷重の強さ $w$ と長さ $L$ の4乗」、分母は「曲げにくさ $EI$」。$wL$ は荷重の合計 $W$ なので、$\delta=\dfrac{5WL^3}{384EI}$ とも書ける。荷重の合計が同じなら、長さの3乗に比例する（①の実験）。384 は、2回の積分と、まん中の値を入れたときに出てくる分数をまとめた数です。
\end{phyomi}
\begin{phbaisuru}{長さ $L$}
本の並べ方（1 m あたりの重さ $w$）が同じなら、$L^4$ で16倍。$L$ が2倍なら本の量も2倍になり、$W$ が2倍、さらに $L^3$ で8倍です。長い棚の中に、支えの板を入れるのはこのためです（例題12）。
\end{phbaisuru}

\bksubsection{9}{代表的なはりの答え}
\label{ph:力-10:sec9}よく出てくる4つのはりの答えを、表にまとめます。どれも⑧と同じ手順（$M$ を書く→2回積分→境界条件）で導けます。片持ちはりの根もとでは「下がらない」と「傾かない」の2つが境界条件です。

\begin{bkcenter}
\begin{bktable}{colspec={Q[c,m,wd=50mm] Q[c,m,wd=24mm] Q[c,m,wd=24mm] Q[c,m,wd=20mm]}, rowsep=2.5pt}
はりと荷重 & いちばん大きい $M$ & たわみ $\delta$ & ばね定数 $k$ \\
{\begin{tikzpicture}[x=1mm, y=1mm]\pxkwall{(0,1)}{8}\fill[bkOrange!35] (0,0) rectangle (32,2); \draw[bkOrange!70!black, line width=0.5pt] (0,0) rectangle (32,2); \pxkload{(32,2.3)}{6}{$F$}\end{tikzpicture}\\片持ち・先に $F$} & {$FL$\\（根もと）} & $\dfrac{FL^3}{3EI}$ & $\dfrac{3EI}{L^3}$ \\
{\begin{tikzpicture}[x=1mm, y=1mm]\pxkwall{(0,1)}{8}\fill[bkOrange!35] (0,0) rectangle (32,2); \draw[bkOrange!70!black, line width=0.5pt] (0,0) rectangle (32,2); \pxkdist{0}{32}{2}{}\end{tikzpicture}\\片持ち・等分布 $w$} & {$\dfrac{wL^2}{2}$\\（根もと）} & $\dfrac{wL^4}{8EI}$ & ― \\
{\begin{tikzpicture}[x=1mm, y=1mm]\fill[bkOrange!35] (0,0) rectangle (36,2); \draw[bkOrange!70!black, line width=0.5pt] (0,0) rectangle (36,2); \pxkpin{(1,0)} \pxkroller{(35,0)} \pxkload{(18,2.3)}{6}{$F$}\end{tikzpicture}\\単純支持・中央に $F$} & {$\dfrac{FL}{4}$\\（中央）} & $\dfrac{FL^3}{48EI}$ & $\dfrac{48EI}{L^3}$ \\
{\begin{tikzpicture}[x=1mm, y=1mm]\fill[bkOrange!35] (0,0) rectangle (36,2); \draw[bkOrange!70!black, line width=0.5pt] (0,0) rectangle (36,2); \pxkpin{(1,0)} \pxkroller{(35,0)} \pxkdist{0}{36}{2}{}\end{tikzpicture}\\単純支持・等分布 $w$} & {$\dfrac{wL^2}{8}$\\（中央）} & $\dfrac{5wL^4}{384EI}$ & ― \\
\end{bktable}
\bktabcaption{代表的なはりの答え（$L$ ははりの長さ。単純支持は「片方はピン、もう片方はころ」で、はしは回れる）}
\end{bkcenter}

\pxkhead{片持ちはりで確かめる}
\label{ph:力-10:katamochi}かべから $x$ の所で切り、右の部分を考えます。右の部分の荷重 $F$ は、切り口を、うで $(L-x)$ で、山の形（上に凸）に曲げます。山の形はマイナスと決めたので $M=-F(L-x)$。
\[ EI\frac{\dd^2y}{\dd x^2}=F(L-x),\qquad EI\frac{\dd y}{\dd x}=F\Bigl(Lx-\frac{x^2}{2}\Bigr)+C_1 \]
\[ EIy=F\Bigl(\frac{Lx^2}{2}-\frac{x^3}{6}\Bigr)+C_1x+C_2 \]
根もと（$x=0$）で傾かない・下がらないので $C_1=C_2=0$。先（$x=L$）では $\delta=\dfrac{F}{EI}\Bigl(\dfrac{L^3}{2}-\dfrac{L^3}{6}\Bigr)=\dfrac{FL^3}{3EI}$ です。

\begin{phbaisuru}{長さ $L$}
\par\vspace{3pt}\noindent\begin{minipage}[c]{0.46\linewidth}
集中荷重の $\delta=\dfrac{FL^3}{48EI}$ では $L^3$ なので8倍（①の実験）。逆に、高さ $h$ を2倍にすると $I$ が8倍で、$\delta$ は $\dfrac{1}{8}$。はば $b$ を2倍にすると $\dfrac{1}{2}$。ヤング率 $E$ が2倍でも $\dfrac{1}{2}$。
\end{minipage}\hfill
\begin{minipage}[c]{0.50\linewidth}
\centering
\begin{tikzpicture}[x=1mm, y=1mm]
\begin{axis}[bkplot, width=62mm, height=42mm, xmin=0, xmax=2.3, ymin=0, ymax=9, xtick={0.5,1,1.5,2}, ytick={1,2,4,8}, xlabel={倍}, ylabel={$\delta$ の倍率}]
  \addplot[bkRed, line width=1.6pt, domain=0:2.08, samples=60] {x^3};
  \addplot[bkBlue, line width=1.6pt, domain=0.48:2.2, samples=60] {1/x^3};
  \addplot[bkOrange, line width=1.6pt, domain=0.12:2.2, samples=60] {1/x};
  \node[anchor=west, font=\sffamily\footnotesize, text=bkRed!80!black] at (axis cs:1.55,7.6) {長さ $L$};
  \node[anchor=west, font=\sffamily\footnotesize, text=bkBlue!80!black] at (axis cs:0.75,6.2) {高さ $h$};
  \node[anchor=west, font=\sffamily\footnotesize, text=bkOrange!80!black] at (axis cs:0.25,5.2) {はば $b$};
\end{axis}
\end{tikzpicture}
\end{minipage}
\end{phbaisuru}

\pxkhead{板もばねになる}
単純支持のはりのまん中に $F$ をのせると、$\delta=\dfrac{FL^3}{48EI}$ です。$k=\dfrac{48EI}{L^3}$ とおくと $F=k\delta$。たわみは力に比例し、はりは「ばね定数 $k$ のばね」です。表の右の列が、そのばね定数です。自動車の板ばね（何枚も重ねた鋼の板）は、この性質を使っています。

\Needspace{14\baselineskip}
\begin{phflow}
\phflowrow{$EI\dfrac{\dd^2y}{\dd x^2}=-M$}{まん中に $F$、両はし\\支え、たわみは小さい}{$F=k\delta$\\（$k=\dfrac{48EI}{L^3}$）}
\phflowrow{$F=k\delta$}{板ばねを\\コイルばねに}{$F=kx$\\（ばねののび、力-8）}
\end{phflow}

\bksubsection{10}{重ね合わせと、不静定のはり}
\label{ph:力-10:sec10}たわみの微分方程式 $EI\dfrac{\dd^2y}{\dd x^2}=-M$ は、$y$ について1次の式です（$y$ を2倍にすると、左辺もちょうど2倍）。2つの荷重による $M$ は、それぞれの $M$ の和なので、たわみも\textbf{それぞれのたわみの和}になります。これを\textbf{重ね合わせ}といいます。

\begin{bkfigure}
\begin{tikzpicture}[x=1mm, y=1mm]
  \begin{scope}[shift={(0,0)}]
    \fill[bkOrange!35] (0,0) rectangle (32,2); \draw[bkOrange!70!black, line width=0.5pt] (0,0) rectangle (32,2);
    \pxkpin{(1,0)} \pxkroller{(31,0)} \pxkdist{0}{32}{2}{} \pxkload{(16,7.4)}{6}{$F$}
    \node[anchor=north, font=\sffamily\footnotesize] at (16,-7) {本＋置物};
  \end{scope}
  \node[font=\sffamily\large] at (38,1) {＝};
  \begin{scope}[shift={(44,0)}]
    \fill[bkOrange!35] (0,0) rectangle (32,2); \draw[bkOrange!70!black, line width=0.5pt] (0,0) rectangle (32,2);
    \pxkpin{(1,0)} \pxkroller{(31,0)} \pxkdist{0}{32}{2}{}
    \node[anchor=north, font=\sffamily\footnotesize] at (16,-7) {$\dfrac{5wL^4}{384EI}$};
  \end{scope}
  \node[font=\sffamily\large] at (82,1) {＋};
  \begin{scope}[shift={(88,0)}]
    \fill[bkOrange!35] (0,0) rectangle (32,2); \draw[bkOrange!70!black, line width=0.5pt] (0,0) rectangle (32,2);
    \pxkpin{(1,0)} \pxkroller{(31,0)} \pxkload{(16,2.3)}{6}{$F$}
    \node[anchor=north, font=\sffamily\footnotesize] at (16,-7) {$\dfrac{FL^3}{48EI}$};
  \end{scope}
\end{tikzpicture}
\bkfigcaption{荷重を分けて、それぞれのたわみを足す（重ね合わせ）}
\end{bkfigure}

\pxkhead{不静定のはり：つり合いの式だけでは反力が決まらない}
3つの支点で支えた長い棚板を考えます。つり合いの式は「力の和＝0」「モーメントの和＝0」の2つですが、知りたい反力は3つです。式が足りないので、つり合いだけでは反力が決まりません。このようなはりを\textbf{不静定のはり}といいます（力-5 の不静定）。

そこで、「まん中の支点では下がらない」という\textbf{たわみの条件}を使います。重ね合わせで考えると、
\begin{enumerate}[label={\bkcircled{\arabic*}}, leftmargin=2.2\zw, itemsep=1pt, topsep=2pt]
\item まん中の支点をはずした長さ $2L$ のはりに $w$ がのると、まん中は $\dfrac{5w(2L)^4}{384EI}$ 下がる。
\item まん中に上向きの力 $R$ だけをかけると、まん中は $\dfrac{R(2L)^3}{48EI}$ 上がる。
\item 本当は下がらないので、①と②が等しい：$\dfrac{5w\cdot16L^4}{384EI}=\dfrac{R\cdot8L^3}{48EI}$
\end{enumerate}
\[ R=\frac{5\times16\times48}{384\times8}wL=\frac{5}{4}wL=1.25\,wL \]
両はしの反力は、残りを半分ずつで $\dfrac{2wL-1.25wL}{2}=0.375\,wL$。まん中の支点は、両はしの3倍以上の力を受けます。

\begin{motto}[5]
\phboxhead{両はしを固定したはり}
両はしがかべにうめこまれたはり（両端固定）は、両はしで「下がらない」「傾かない」の4つの条件があり、これも不静定です。はしの曲げモーメントを $-M_0$（山の形）とおくと、等分布荷重では $M=\dfrac{w}{2}(Lx-x^2)-M_0$。⑧と同じように2回積分して、4つの条件から $C_1$、$C_2$、$M_0$ を決めると、$M_0=\dfrac{wL^2}{12}$、まん中のたわみは $\delta=\dfrac{wL^4}{384EI}$ で、単純支持の $\dfrac{1}{5}$ になります（例題13）。建物のはりと柱をしっかりつなぐと、はりのたわみが小さくなるのはこのためです。
\end{motto}

\begin{seikaku}
\phboxhead{2つの近似}
曲がり方の正確な式は $\dfrac{1}{\rho}=\dfrac{|y''|}{\bigl(1+(y')^2\bigr)^{3/2}}$ です（力-1 ⑤）。本棚の板の傾きは最大でも約0.009なので、$(y')^2$ は約0.00008で、1にくらべて無視できます。また、「断面は平らなまま」という仮定は、長さが高さの10倍以上ある細長いはりではよく合いますが、短く太いはりでは、せん断力によるゆがみが加わって、たわみがこの章の式より大きくなります。
\end{seikaku}

\begin{phmada}[骨は、力に合わせて形を変える]
動物の骨の中には、細い骨の柱（こつりょう）がスポンジのように組み合わさっています。この細い柱は、骨にかかる力の向きにそって並んでいて、使い方が変わると何か月もかけて並び方を変えます。1892年にドイツの医師ウォルフがこの考えをまとめ、「ウォルフの法則」とよばれています。骨の中の細胞（骨細胞）が、ひずみや、骨の中の液体の流れを感じて、骨を作る細胞とこわす細胞に合図を出していると考えられています。でも、細胞がどんなしくみで力を感じ、どれだけの力でどう作りかえるのかは、まだ一部しかわかっておらず、研究が続いています。
\end{phmada}

\begin{phsaki}{板・殻・振動・コンピュータ}
この章のはりは「細長い棒」でした。床のように広い板や、タンクのように曲がった板（殻）の曲げは、2つの向きに同時に曲がるので、偏微分方程式（数学 8-13）になります。はりがゆれるときのはりの振動は、$EI\dfrac{\partial^4y}{\partial x^4}$ に慣性の項が加わった式で、波-6 で学びます。形が複雑な構造は、小さな要素に分けてコンピュータで解きます（有限要素法、力-16）。これらは大学の材料力学・構造力学の上の学年で学びます。
\end{phsaki}

\begin{senmonsho}[大学の教科書ふうの書き方]
\begin{bkhonbun}
はりの微小要素の力とモーメントのつり合いより $\dfrac{\dd F}{\dd x}=-q$、$\dfrac{\dd M}{\dd x}=F$ を得る。平面保持の仮定のもとで曲げ応力は $\sigma=\dfrac{M}{I}y$ となり、たわみ曲線の微分方程式は $EI\dfrac{\dd^2v}{\dd x^2}=-M$ で与えられる。
\end{bkhonbun}
\bkyaku{専門書によって、せん断力を $F$ や $V$、分布荷重を $q$、たわみを $v$ や $w$ と書く。この章の $Q$・$w$・$y$ と同じもの。どの文字を使うかは本によってちがうので、はじめの図で確かめよう。}
\end{senmonsho}

\bkkigocards{\bkkigocard{$I$}{アイ}{断面二次モーメント。曲げにくさを決める断面の形の量（m$^4$）。力-6 の慣性モーメント $I$ とは別の量}{second moment of area}{力-10}}{\bkkigocard{$EI$}{イーアイ}{曲げ剛性。材料のかたさ $E$ と断面の $I$ のかけ算（N·m$^2$）}{flexural rigidity}{力-10}}

\begin{kotoba}\small
\bkword{はり}{細長く、横から力を受けて曲がる部材。}
\bkword{せん断力 $Q$}{切り口で、となりの部分が上下におす力。}
\bkword{曲げモーメント $M$}{切り口で、となりの部分が回そうとする働き。}
\bkword{中立軸}{曲げても、のびもちぢみもしない所。断面の重心を通る。}
\bkword{断面二次モーメント $I$}{$\int y^2\,\dd A$。曲げにくさを決める断面の形の量。}
\bkword{断面係数 $Z$}{$\dfrac{I}{y_{\max}}$。いちばん大きい曲げ応力は $\dfrac{M}{Z}$。}
\bkword{境界条件}{両はしで決まっている条件（下がらない・傾かない）。}
\bkword{不静定}{つり合いの式だけでは、反力が決まらないこと。}
\end{kotoba}

\begin{taisetsu}
\def\tsl#1{\makebox[7.6\zw][l]{\sffamily\bfseries\small\color{\bkTT}#1}}
\tsl{つり合い}$\dfrac{\dd Q}{\dd x}=-w$、\ $\dfrac{\dd M}{\dd x}=Q$（$Q=0$ の所で $M$ が最大）\\[3pt]
\tsl{曲げ応力}$\sigma=\dfrac{My}{I}$、いちばん大きいのは $\sigma_{\max}=\dfrac{M}{Z}$\\[3pt]
\tsl{断面の量}$I=\displaystyle\int y^2\,\dd A$。長方形は $I=\dfrac{bh^3}{12}$、$Z=\dfrac{bh^2}{6}$\\[3pt]
\tsl{たわみ}$EI\dfrac{\dd^2y}{\dd x^2}=-M$。2回積分して、両はしの条件（境界条件）で $C$ を決める\\[3pt]
\tsl{代表的な答え}単純支持の等分布 $\dfrac{5wL^4}{384EI}$、中央の集中 $\dfrac{FL^3}{48EI}$、片持ちの先 $\dfrac{FL^3}{3EI}$\\[3pt]
\tsl{2倍にすると}長さ $L$ は8倍（$L^3$）、高さ $h$ は $\dfrac{1}{8}$、はば $b$ は $\dfrac{1}{2}$\\[3pt]
\tsl{重ね合わせ}荷重ごとのたわみを足してよい。不静定のはりは、たわみの条件を足して解く
\end{taisetsu}

%% ================================================================
\bksection{reidai}{例題でたしかめよう}

\begin{phreidai}[本棚の板のたわみ（はてな？）]
\label{ph:力-10:ex1}長さ\qty{0.8}{m}、はば\qty{0.25}{m}、厚さ\qty{18}{mm}のすぎの板（$E=\qty{7}{GPa}$）を、両はしで本棚の側板にのせました。本\qty{30}{kg}を一様に並べたときの、まん中のたわみを求めましょう。また、同じ本を全部まん中に積んだら、たわみはいくつでしょう。
\phbamen
板は両はしのダボ（小さな棒）の上にのっているだけで、はしは自由に回れます。知りたいのは、まん中が下がる量です。
\phmushi
板の自分の重さは、本より小さいので無視します（\qty{0.8}{m}のすぎの板は約\qty{1.4}{kg}）。本の重さは板の上に一様にかかるとします。たわみは小さく、板は細長いとします。
\phzu
\par\vspace{3pt}\noindent\begin{minipage}[c]{0.50\linewidth}
\begin{tikzpicture}[x=1mm, y=1mm]
  \fill[bkOrange!35] (0,0) rectangle (52,2.4); \draw[bkOrange!70!black, line width=0.5pt] (0,0) rectangle (52,2.4);
  \pxkpin{(0,0)} \pxkroller{(52,0)} \pxkdist{0}{52}{2.4}{$w$}
  \draw[bkGray, line width=0.4pt, <->] (0,-9) -- (52,-9); \node[anchor=north, font=\sffamily\footnotesize] at (26,-9.4) {$L=\qty{0.8}{m}$};
\end{tikzpicture}
\end{minipage}\hfill
\begin{minipage}[c]{0.46\linewidth}
はりの長さ $L$、分布荷重 $w$（N/m）、ヤング率 $E$、断面二次モーメント $I$、まん中のたわみ $\delta$。
\end{minipage}
\phriyu
両はしで支えた（単純支持）はりに、等分布の荷重がのっているので、⑧で導いた $\delta=\dfrac{5wL^4}{384EI}$ を使います。まん中に積むときは、⑨の表の $\delta=\dfrac{FL^3}{48EI}$ です。
\phtateru
\[ w=\frac{30\times9.8}{0.8}=\pxkq{367.5}{N/m},\qquad I=\frac{bh^3}{12}=\frac{0.25\times0.018^3}{12} \]
\phtoku
\[ I=\qty{1.215e-7}{m^4},\qquad EI=7\times10^9\times1.215\times10^{-7}=\pxkq{850.5}{N.m^2} \]
\[ \delta=\frac{5\times367.5\times0.8^4}{384\times850.5}=\frac{752.6}{\num{326592}}=\qty{2.30e-3}{m} \]
まん中に積むと $F=\qty{294}{N}$ で、$\delta=\dfrac{294\times0.8^3}{48\times850.5}=\qty{3.69e-3}{m}$。
\phtashikame
単位：$\dfrac{\pxku{N/m}\times\unit{m^4}}{\pxku{N.m^2}}=\unit{m}$ で長さ。\\
大きさ：\qty{2.3}{mm}は、目で見てわかるくらいの小さなたわみです。本棚の板として、ありそうな値です。\\
極端な場合：$E$ や $I$ を大きく（かたく）すると $\delta$ は0に近づきます。まん中に積むと、たわみは $\dfrac{3.69}{2.30}=1.6$ 倍。重い物は、はしに置くほうがよいことがわかります。
\kotae 約\qty{2.3}{mm}。まん中に積むと約\qty{3.7}{mm}
\end{phreidai}
\begin{jibun}
長さ\qty{0.9}{m}、はば\qty{0.30}{m}、厚さ\qty{20}{mm}のすぎの板（$E=\qty{7}{GPa}$）を両はしで支え、本\qty{40}{kg}を一様に並べました。まん中のたわみを求めましょう。
\bkjans{$w=\pxkq{435.6}{N/m}$、$I=\qty{2.0e-7}{m^4}$、$EI=\pxkq{1400}{N.m^2}$。$\delta=\frac{5wL^4}{384EI}\approx\qty{2.7}{mm}$}
\end{jibun}

\begin{phreidai}[定規を寝かせるか、立てるか]
\label{ph:力-10:ex2}はば\qty{30}{mm}、厚さ\qty{3}{mm}の定規を、両はしで支えます。寝かせたときと、立てたときで、同じ力をまん中にかけたときのたわみは何倍ちがうでしょう。
\phbamen
同じ定規・同じ長さ・同じ力で、向きだけを変えます。
\phmushi
定規の自分の重さと、目もりのみぞは無視し、断面は長方形とします。
\phzu
寝かせたとき：$b=\qty{30}{mm}$、$h=\qty{3}{mm}$。立てたとき：$b=\qty{3}{mm}$、$h=\qty{30}{mm}$（$h$ は力の向きに測った長さ）。
\phriyu
たわみ $\delta=\dfrac{FL^3}{48EI}$ で、$F$・$L$・$E$ は同じなので、$\delta$ は $I$ に反比例します。$I=\dfrac{bh^3}{12}$ をくらべます。
\phtateru
\[ I_{\text{寝}}=\frac{30\times3^3}{12},\qquad I_{\text{立}}=\frac{3\times30^3}{12} \]
\phtoku
\[ I_{\text{寝}}=\qty{67.5}{mm^4},\qquad I_{\text{立}}=\qty{6750}{mm^4},\qquad \frac{I_{\text{立}}}{I_{\text{寝}}}=100 \]
\phtashikame
単位：mm$^4$ どうしの比なので、単位はありません。\\
別の考え方で：$\dfrac{I_{\text{立}}}{I_{\text{寝}}}=\dfrac{3\times30^3}{30\times3^3}=\Bigl(\dfrac{30}{3}\Bigr)^2=100$。高さとはばの比の2乗です。\\
極端な場合：はばと厚さが同じ（正方形）なら、比は1で、どちら向きでも同じです。
\kotae 立てると、たわみは寝かせたときの $\dfrac{1}{100}$
\end{phreidai}
\begin{jibun}
はば\qty{25}{mm}、厚さ\qty{2}{mm}のものさしを、寝かせたときと立てたときで、同じ力によるたわみは何倍ちがうでしょう。
\bkjans{$(25\div2)^2=156.25$ で、立てると約 $\frac{1}{156}$}
\end{jibun}

\begin{phreidai}[2点で支えた板が受ける力と、いちばん大きい曲げモーメント]
\label{ph:力-10:ex3}長さ\qty{2.0}{m}の板を両はしA・Bで支え、Aから\qty{0.5}{m}の所に、重さ\qty{600}{N}の荷物をのせました。A・Bが板を支える力と、いちばん大きい曲げモーメントを求め、$Q$ 図と $M$ 図をかきましょう。
\phbamen
みぞの上にわたした板の上に、荷物を置いた場面です。荷物はAに近い所にあります。
\phmushi
板の重さは無視します。荷物は1点にのっているとします。
\phzu
\par\vspace{3pt}\noindent\begin{minipage}[c]{0.50\linewidth}
\begin{tikzpicture}[x=1mm, y=1mm]
  \fill[bkOrange!35] (0,0) rectangle (52,2.4); \draw[bkOrange!70!black, line width=0.5pt] (0,0) rectangle (52,2.4);
  \pxkpin{(0,0)} \pxkroller{(52,0)} \pxkload{(13,2.7)}{7}{$F=\qty{600}{N}$}
  \draw[bkGray, line width=0.4pt, <->] (0,-9) -- (13,-9); \node[anchor=north, font=\sffamily\footnotesize] at (6.5,-9.4) {$a=\qty{0.5}{m}$};
  \draw[bkGray, line width=0.4pt, <->] (13,-9) -- (52,-9); \node[anchor=north, font=\sffamily\footnotesize] at (32.5,-9.4) {$b=\qty{1.5}{m}$};
  \node[anchor=east, font=\sffamily\footnotesize] at (-1.5,1.2) {A}; \node[anchor=west, font=\sffamily\footnotesize] at (53.5,1.2) {B};
\end{tikzpicture}
\end{minipage}\hfill
\begin{minipage}[c]{0.46\linewidth}
支点の力 $R_{\mathrm A}$・$R_{\mathrm B}$（上向き）、荷物までの長さ $a$・$b$、板の長さ $L$。
\end{minipage}
\phriyu
荷物が1つだけなので、⑥で導いた $R_{\mathrm A}=\dfrac{b}{L}F$、$R_{\mathrm B}=\dfrac{a}{L}F$ を使います。$M$ は $Q$ を足したものなので、$Q$ が正から負へ変わる荷物の所で最大です。
\phtateru
\[ R_{\mathrm A}=\frac{1.5}{2.0}\times600,\qquad R_{\mathrm B}=\frac{0.5}{2.0}\times600,\qquad M_{\max}=R_{\mathrm A}\times a \]
\phtoku
\[ R_{\mathrm A}=\qty{450}{N},\qquad R_{\mathrm B}=\qty{150}{N},\qquad M_{\max}=450\times0.5=\pxkq{225}{N.m} \]
$Q$ 図：A〜荷物は\qty{450}{N}、荷物〜Bは $450-600=-\qty{150}{N}$。$M$ 図：Aで0から、荷物の所の\pxkq{225}{N.m}まで直線で上がり、Bで0まで直線で下がる三角形です。
\begin{center}
\begin{tikzpicture}[x=1mm, y=1mm]
\begin{axis}[bkplot, axis lines=left, width=52mm, height=34mm, xmin=0, xmax=2.2, ymin=-200, ymax=520, xtick={0.5,1,1.5,2}, ytick={-150,0,450}, xlabel={$x$}, ylabel={$Q$（N）}]
  \draw[bkInk, line width=0.6pt] (axis cs:0,0) -- (axis cs:2.2,0);
  \fill[bkBlue!15] (axis cs:0,0) -- (axis cs:0,450) -- (axis cs:0.5,450) -- (axis cs:0.5,-150) -- (axis cs:2,-150) -- (axis cs:2,0) -- cycle;
  \draw[bkBlue, line width=1.5pt] (axis cs:0,0) -- (axis cs:0,450) -- (axis cs:0.5,450) -- (axis cs:0.5,-150) -- (axis cs:2,-150) -- (axis cs:2,0);
\end{axis}
\begin{axis}[bkplot, at={(60mm,0)}, width=52mm, height=34mm, xmin=0, xmax=2.2, ymin=0, ymax=270, xtick={0.5,1,1.5,2}, ytick={225}, xlabel={$x$}, ylabel={$M$（N·m）}]
  \fill[bkRed!12] (axis cs:0,0) -- (axis cs:0.5,225) -- (axis cs:2,0) -- cycle;
  \draw[bkRed, line width=1.5pt] (axis cs:0,0) -- (axis cs:0.5,225) -- (axis cs:2,0);
\end{axis}
\end{tikzpicture}
\end{center}
\phtashikame
単位：N と N·m。\\
力の和：$450+150=600$ で、荷物の重さと等しい。\\
別の方法で：Bのまわりでも $M$ を計算すると $R_{\mathrm B}\times b=150\times1.5=\pxkq{225}{N.m}$。同じです。\\
極端な場合：荷物がまん中（$a=b=\qty{1}{m}$）なら、$R_{\mathrm A}=R_{\mathrm B}=\qty{300}{N}$、$M_{\max}=\dfrac{FL}{4}=\pxkq{300}{N.m}$ で最大。荷物がAの真上なら、$R_{\mathrm A}=\qty{600}{N}$ で、$M$ は0（板は曲がらない）。
\kotae $R_{\mathrm A}=\qty{450}{N}$、$R_{\mathrm B}=\qty{150}{N}$、$M_{\max}=\pxkq{225}{N.m}$（荷物の所）
\end{phreidai}
\begin{jibun}
長さ\qty{3.0}{m}の板を両はしA・Bで支え、Aから\qty{1.0}{m}の所に\qty{900}{N}の荷物をのせました。A・Bの支える力と、いちばん大きい曲げモーメントを求めましょう。
\bkjans{$R_{\mathrm A}=\qty{600}{N}$、$R_{\mathrm B}=\qty{300}{N}$、$M_{\max}=600\times1.0=\pxkq{600}{N.m}$}
\end{jibun}

\begin{phreidai}[はみ出した板の上を歩く：板が浮き上がるのはどこ？]
\label{ph:力-10:ex4}長さ\qty{4}{m}、質量\qty{20}{kg}の一様な板を、左はしAと、Aから\qty{3}{m}の点Bの2か所で支えました（右の\qty{1}{m}は、Bの外へはみ出している）。質量\qty{60}{kg}の人が、Aから右へ歩いていきます。人がBの外へどこまで進むと、Aの所で板が浮き上がり始めるでしょう。
\phbamen
ベランダや工事現場の足場で、支えの外へはみ出した板（はね出し）に乗る場面です。板はAの上にのっているだけで、くぎでとめていません。
\phmushi
板はたわまない（まっすぐな棒）とし、人は1点に体重をかけるとします。
\phzu
\par\vspace{3pt}\noindent\begin{minipage}[c]{0.52\linewidth}
\begin{tikzpicture}[x=1mm, y=1mm]
  \fill[bkOrange!35] (0,0) rectangle (56,2.4); \draw[bkOrange!70!black, line width=0.5pt] (0,0) rectangle (56,2.4);
  \pxkpin{(0,0)} \pxkpin{(42,0)}
  \draw[bk arrow, bkGray, line width=1.1pt] (28,2.4) -- (28,-7); \node[anchor=north, font=\sffamily\footnotesize] at (28,-7.3) {\qty{196}{N}};
  \bkperson{(49,2.4)}{13}{bkBlue}
  \draw[bk arrow, bkBlue, line width=1.1pt] (49,2.4) -- (49,-7); \node[anchor=north, font=\sffamily\footnotesize] at (49.5,-7.3) {\qty{588}{N}};
  \node[anchor=south, font=\sffamily\footnotesize] at (0,3) {A}; \node[anchor=south east, font=\sffamily\footnotesize] at (41,3) {B};
  \draw[bkGray, line width=0.4pt, <->] (0,-15) -- (42,-15); \node[anchor=north, font=\sffamily\footnotesize] at (21,-15.4) {\qty{3}{m}};
  \draw[bkGray, line width=0.4pt, <->] (42,-15) -- (49,-15); \node[anchor=north, font=\sffamily\footnotesize] at (45.5,-15.4) {$x$};
\end{tikzpicture}
\end{minipage}\hfill
\begin{minipage}[c]{0.44\linewidth}
板の重さ $20\times9.8=\qty{196}{N}$ は、板のまん中（Aから\qty{2}{m}、Bの左\qty{1}{m}）にかかる。人の重さ $60\times9.8=\qty{588}{N}$。人はBから右へ $x$ の所。
\end{minipage}
\phriyu
板が浮き上がり始めるのは、Aが板を支える力 $R_{\mathrm A}$ が0になるときです（Aは板を引っぱれない）。$R_{\mathrm A}$ を求めるには、Bのまわりのモーメントのつり合い（てこのつり合い）を使います。
\phtateru
Bのまわりで、左回り（板の重さ）と右回り（人）の回す働きを考えます。$R_{\mathrm A}=0$ のときは、
\[ 196\times1=588\times x \]
\phtoku
\[ x=\frac{196}{588}=\frac{1}{3}\approx\qty{0.33}{m} \]
人がBの真上（$x=0$）にいるときは、Bのまわりで $R_{\mathrm A}\times3=196\times1$ より $R_{\mathrm A}=\qty{65.3}{N}$。Bの外へ進むほど $R_{\mathrm A}$ はへり、$x=\qty{0.33}{m}$ で0になります。
\phtashikame
単位：$\dfrac{\unit{N}\times\unit{m}}{\unit{N}}=\unit{m}$。\\
大きさ：はみ出しは\qty{1}{m}ありますが、\qty{0.33}{m}先までしか進めません。板の重さが人より軽いからです。\\
極端な場合：板がとても重ければ、はしまで歩けます（$196\times1\geqq588\times1$ なら浮かない）。人がとても重ければ、Bをこえたとたんに浮きます。
\kotae Bから約\qty{0.33}{m}（Aから約\qty{3.3}{m}）
\end{phreidai}
\begin{jibun}
長さ\qty{5}{m}、質量\qty{30}{kg}の一様な板を、左はしAと、Aから\qty{4}{m}の点Bで支えました。質量\qty{50}{kg}の人がBの外へ歩くとき、Bからどこまで進むと、Aの所で板が浮き上がり始めるでしょう。
\bkjans{板の重さはBの左\qty{1.5}{m}。$30\times9.8\times1.5=50\times9.8\times x$ より $x=\qty{0.9}{m}$}
\end{jibun}

\begin{phreidai}[本棚の板のせん断力・曲げモーメント・曲げ応力]
\label{ph:力-10:ex5}例題1の本棚の板（長さ\qty{0.8}{m}、はば\qty{0.25}{m}、厚さ\qty{18}{mm}、本\qty{30}{kg}を一様に）で、いちばん大きいせん断力・曲げモーメント・曲げ応力を求めましょう。
\phbamen
板がこわれないかを確かめます。曲げ応力は、板の上の面と下の面で最大です。
\phmushi
板の自分の重さは無視し、本の重さは一様にかかるとします。
\phzu
例題1と同じ図。$w=\pxkq{367.5}{N/m}$、断面ははば $b$、高さ $h$ の長方形。
\phriyu
等分布の単純支持なので、$Q$ は両はしで最大 $\dfrac{wL}{2}$、$M$ はまん中で最大 $\dfrac{wL^2}{8}$（④の表、⑨の表）。曲げ応力の最大は $\sigma_{\max}=\dfrac{M}{Z}$、$Z=\dfrac{bh^2}{6}$（⑦）。
\phtateru
\[ Q_{\max}=\frac{367.5\times0.8}{2},\quad M_{\max}=\frac{367.5\times0.8^2}{8},\quad Z=\frac{0.25\times0.018^2}{6} \]
\phtoku
\[ Q_{\max}=\qty{147}{N},\quad M_{\max}=\pxkq{29.4}{N.m},\quad Z=\qty{1.35e-5}{m^3} \]
\[ \sigma_{\max}=\frac{29.4}{1.35\times10^{-5}}=\qty{2.18e6}{Pa}\approx\qty{2.2}{MPa} \]
\phtashikame
単位：$\dfrac{\pxku{N.m}}{\unit{m^3}}=\unit{N/m^2}=\unit{Pa}$。\\
大きさ：すぎの曲げの強さは数十 MPa あるので、\qty{2.2}{MPa}なら十分にたえます。本棚の板は、強さよりも「たわみ」で厚さが決まることが多いのです。\\
別の方法で：③の表計算の $M$ の最大も\pxkq{29.4}{N.m}でした。
\kotae $Q_{\max}=\qty{147}{N}$（両はし）、$M_{\max}=\pxkq{29.4}{N.m}$（まん中）、$\sigma_{\max}\approx\qty{2.2}{MPa}$
\end{phreidai}
\begin{jibun}
じぶんでとこう（例題1）の板（長さ\qty{0.9}{m}、はば\qty{0.30}{m}、厚さ\qty{20}{mm}、本\qty{40}{kg}を一様に）で、いちばん大きいせん断力・曲げモーメント・曲げ応力を求めましょう。
\bkjans{$Q_{\max}=\qty{196}{N}$、$M_{\max}=\pxkq{44.1}{N.m}$、$Z=\qty{2.0e-5}{m^3}$、$\sigma_{\max}\approx\qty{2.2}{MPa}$}
\end{jibun}

\begin{phreidai}[鋼の片持ちはりの先のたわみ]
\label{ph:力-10:ex6}かべから\qty{1.0}{m}つき出した鋼のはり（$E=\qty{205}{GPa}$、$I=\qty{1.0e-6}{m^4}$）の先に、\qty{1.0}{kN}の看板をつるしました。先のたわみと傾きを求めましょう。
\phbamen
かべにしっかり固定した腕の先に、看板がぶら下がっています。
\phmushi
はりの自分の重さは無視します。看板の重さは先の1点にかかるとします。
\phzu
\par\vspace{3pt}\noindent\begin{minipage}[c]{0.46\linewidth}
\begin{tikzpicture}[x=1mm, y=1mm]
  \pxkwall{(0,1.2)}{12}
  \fill[bkGray!40] (0,0) rectangle (44,2.4); \draw[bkInk, line width=0.5pt] (0,0) rectangle (44,2.4);
  \draw[bkInk, line width=0.6pt] (44,0) -- (44,-9);
  \fill[bkYellow!50] (39,-15) rectangle (49,-9); \draw[bkInk, line width=0.5pt] (39,-15) rectangle (49,-9);
  \draw[bkGray, line width=0.4pt, <->] (0,6) -- (44,6); \node[anchor=south, font=\sffamily\footnotesize] at (22,6.3) {$L=\qty{1.0}{m}$};
  \node[anchor=west, font=\sffamily\footnotesize] at (50,-12) {$F$};
\end{tikzpicture}
\end{minipage}\hfill
\begin{minipage}[c]{0.50\linewidth}
はりの長さ $L$、先の荷重 $F$、たわみ $\delta$、先の傾き $\theta$。
\end{minipage}
\phriyu
片持ちはりの先に集中荷重なので、⑨で導いた $\delta=\dfrac{FL^3}{3EI}$ を使います。傾きは、その途中の式 $EI\dfrac{\dd y}{\dd x}=F\Bigl(Lx-\dfrac{x^2}{2}\Bigr)$ に $x=L$ を入れた $\theta=\dfrac{FL^2}{2EI}$ です。
\phtateru
\[ EI=205\times10^9\times1.0\times10^{-6}=\pxkq{2.05e5}{N.m^2} \]
\phtoku
\[ \delta=\frac{1000\times1.0^3}{3\times2.05\times10^5}=\qty{1.63e-3}{m} \]
\[ \theta=\frac{1000\times1.0^2}{2\times2.05\times10^5}=\qty{2.4e-3}{rad} \]
\phtashikame
単位：$\dfrac{\unit{N}\times\unit{m^3}}{\pxku{N.m^2}}=\unit{m}$、傾きは単位なし（rad）。\\
大きさ：\qty{1.6}{mm}は、\qty{1}{m}の腕の $\dfrac{1}{600}$ くらい。\\
極端な場合：長さを2倍にすると $2^3=8$ 倍の\qty{13}{mm}になります。
\kotae たわみ 約\qty{1.6}{mm}、傾き 約\qty{2.4e-3}{rad}（約0.14°）
\end{phreidai}
\begin{jibun}
長さ\qty{1.5}{m}の鋼の片持ちはり（$E=\qty{205}{GPa}$、$I=\qty{2.0e-6}{m^4}$）の先に\qty{500}{N}をかけました。先のたわみを求めましょう。
\bkjans{$\delta=\frac{FL^3}{3EI}\approx\qty{1.4}{mm}$}
\end{jibun}

\begin{phreidai}[2×4材を立てて使う理由]
\label{ph:力-10:ex7}断面が\qty{38}{mm}×\qty{89}{mm}の木材（2×4材）を、長さ\qty{1.8}{m}の単純支持のはりにし、まん中に\qty{1000}{N}をかけます。断面を「立てて（\qty{89}{mm}をたて）」使うときと、「寝かせて（\qty{38}{mm}をたて）」使うときの、いちばん大きい曲げ応力を求めましょう。
\phbamen
日曜大工で作る台や、家の床をつくるはりの場面です。同じ材木でも、向きで強さが変わるでしょうか。
\phmushi
材木の自分の重さは無視し、断面は長方形とします。
\phzu
立てる：$b=\qty{38}{mm}$、$h=\qty{89}{mm}$。寝かせる：$b=\qty{89}{mm}$、$h=\qty{38}{mm}$。$L=\qty{1.8}{m}$、$F=\qty{1000}{N}$。
\phriyu
こわれにくさは曲げ応力で決まります。$\sigma_{\max}=\dfrac{M}{Z}$、中央の集中荷重の $M_{\max}=\dfrac{FL}{4}$（⑨）、$Z=\dfrac{bh^2}{6}$（⑦）。
\phtateru
\[ M_{\max}=\frac{1000\times1.8}{4}=\pxkq{450}{N.m} \]
\[ Z_{\text{立}}=\frac{0.038\times0.089^2}{6},\qquad Z_{\text{寝}}=\frac{0.089\times0.038^2}{6} \]
\phtoku
\[ Z_{\text{立}}=\qty{5.02e-5}{m^3},\qquad Z_{\text{寝}}=\qty{2.14e-5}{m^3} \]
\[ \sigma_{\text{立}}=\frac{450}{5.02\times10^{-5}}\approx\qty{9.0}{MPa},\qquad \sigma_{\text{寝}}=\frac{450}{2.14\times10^{-5}}\approx\qty{21}{MPa} \]
\phtashikame
単位：Pa。\\
比で確かめる：$\dfrac{Z_{\text{立}}}{Z_{\text{寝}}}=\dfrac{89}{38}=2.34$ 倍（強さ）。たわみの比は $I$ の比で $\Bigl(\dfrac{89}{38}\Bigr)^2=5.5$ 倍。強さは高さの2乗、かたさは高さの3乗できいてきます。\\
極端な場合：正方形の断面なら、どちら向きでも同じです。
\kotae 立てる 約\qty{9.0}{MPa}、寝かせる 約\qty{21}{MPa}（立てるほうが約2.3倍強い）
\end{phreidai}
\begin{jibun}
\qty{105}{mm}角（正方形）の木材を長さ\qty{2.0}{m}の単純支持のはりにし、まん中に\qty{2000}{N}をかけました。いちばん大きい曲げ応力を求めましょう。
\bkjans{$M=\pxkq{1000}{N.m}$、$Z=\qty{1.93e-4}{m^3}$、$\sigma\approx\qty{5.2}{MPa}$}
\end{jibun}

\begin{phreidai}[H形鋼は、なぜその形？]
\label{ph:力-10:ex8}高さ\qty{200}{mm}、はば\qty{100}{mm}、まん中の板（ウェブ）の厚さ\qty{5.5}{mm}、上下の板（フランジ）の厚さ\qty{8}{mm}のH形鋼の $I$ を求め、同じ断面積の正方形の棒とくらべましょう。
\phbamen
ビルや橋のはりには、H の形の断面の鋼（H形鋼）がよく使われます。同じ量の鋼を、どんな形に並べると曲がりにくいのでしょう。
\phmushi
本物のH形鋼の、板と板のつなぎ目の丸み（すみの丸み）は無視します。
\phzu
\par\vspace{3pt}\noindent\begin{minipage}[c]{0.36\linewidth}
\centering
\begin{tikzpicture}[x=1mm, y=1mm]
  \fill[bkBlue!25] (0,0) rectangle (20,1.6) (9.45,1.6) rectangle (10.55,38.4) (0,38.4) rectangle (20,40);
  \draw[bkBlue!70!black, line width=0.6pt] (0,0) -- (20,0) -- (20,1.6) -- (10.55,1.6) -- (10.55,38.4) -- (20,38.4) -- (20,40) -- (0,40) -- (0,38.4) -- (9.45,38.4) -- (9.45,1.6) -- (0,1.6) -- cycle;
  \draw[bkRed, dashed, line width=0.5pt] (-3,20) -- (23,20);
  \draw[bkGray, line width=0.4pt, <->] (-5,0) -- (-5,40); \node[anchor=east, font=\sffamily\footnotesize] at (-5.5,30) {200};
  \draw[bkGray, line width=0.4pt, <->] (0,43) -- (20,43); \node[anchor=south, font=\sffamily\footnotesize] at (10,43.3) {100};
  \node[anchor=west, font=\sffamily\footnotesize] at (21,39) {8};
  \node[anchor=west, font=\sffamily\footnotesize] at (12,26) {5.5};
\end{tikzpicture}
\end{minipage}\hfill
\begin{minipage}[c]{0.60\linewidth}
外の長方形（\qty{100}{mm}×\qty{200}{mm}）から、左右の欠けた部分（合わせて はば $100-5.5=\qty{94.5}{mm}$、高さ $200-2\times8=\qty{184}{mm}$ の長方形）を取りのぞいた形と考えます。中立軸は、上下対称なのでまん中（赤い点線）です。単位は mm で計算します。
\end{minipage}
\phriyu
$I=\displaystyle\int y^2\,\dd A$ は足し算なので、「外の長方形の $I$」から「欠けた長方形の $I$」をひけばよい（どちらも中立軸が同じ所）。長方形は $I=\dfrac{bh^3}{12}$。
\phtateru
\[ I_{\mathrm H}=\frac{100\times200^3}{12}-\frac{94.5\times184^3}{12} \]
断面積：$A=2\times100\times8+184\times5.5=\qty{2612}{mm^2}$。同じ面積の正方形の1辺は $\sqrt{2612}=\qty{51.1}{mm}$。
\phtoku
\[ I_{\mathrm H}=\num{66666667}-\num{49057344}=\qty{1.76e7}{mm^4} \]
\[ I_{\text{正方形}}=\frac{51.1^4}{12}=\qty{5.69e5}{mm^4},\qquad \frac{I_{\mathrm H}}{I_{\text{正方形}}}\approx31 \]
\phtashikame
単位：mm$^4$。$\qty{1}{mm^4}=\qty{e-12}{m^4}$ なので $I_{\mathrm H}=\qty{1.76e-5}{m^4}$。\\
大きさ：同じ量の鋼で、約31倍曲がりにくくなります。$I$ の約84\,\%は、上下のフランジ（中立軸から遠い所）がかせいでいます。\\
極端な場合：ウェブを0にすると、上下の板がばらばらになり、②の「重ねただけ」と同じでずれてしまいます。ウェブは、上下の板がずれないようにつなぐ役目をしています。
\kotae $I_{\mathrm H}\approx\qty{1.76e7}{mm^4}$（\qty{1760}{cm^4}）。同じ断面積の正方形の棒の約31倍
\end{phreidai}
\begin{jibun}
高さ\qty{150}{mm}、はば\qty{75}{mm}、ウェブの厚さ\qty{5}{mm}、フランジの厚さ\qty{7}{mm}のH形鋼（すみの丸みは無視）の $I$ を求め、同じ断面積の正方形の棒の何倍かを求めましょう。
\bkjans{$I\approx\qty{6.42e6}{mm^4}$。$A=\qty{1730}{mm^2}$、正方形の $I=A^2\div12\approx\qty{2.49e5}{mm^4}$ で、約26倍}
\end{jibun}

\begin{phreidai}[竹と骨はなぜ中が空っぽ？：パイプと棒]
\label{ph:力-10:ex9}外径\qty{40}{mm}、内径\qty{30}{mm}のパイプと、同じ断面積の中の詰まった丸い棒で、$I$ をくらべましょう。
\phbamen
竹や鳥の骨、自転車のフレームは、中が空いたパイプの形です。同じ重さなら、どちらが曲がりにくいでしょう。
\phmushi
材料は同じで、断面はきれいな円とします。
\phzu
パイプ：外径 $D=\qty{40}{mm}$、内径 $d=\qty{30}{mm}$。棒：直径 $d_{\mathrm s}$。
\phriyu
同じ材料・同じ断面積なら、同じ長さで同じ重さです。曲がりにくさは $I$ でくらべます。円は $I=\dfrac{\pi d^4}{64}$、パイプは $I=\dfrac{\pi(D^4-d^4)}{64}$（⑦の表）。
\phtateru
断面積が等しい：$\dfrac{\pi}{4}d_{\mathrm s}^2=\dfrac{\pi}{4}(D^2-d^2)$ より $d_{\mathrm s}=\sqrt{40^2-30^2}$
\phtoku
\[ d_{\mathrm s}=\sqrt{700}=\qty{26.5}{mm} \]
\[ I_{\text{パイプ}}=\frac{\pi(40^4-30^4)}{64}=\qty{8.59e4}{mm^4} \]
\[ I_{\text{棒}}=\frac{\pi\times26.46^4}{64}=\qty{2.41e4}{mm^4} \]
\[ \frac{I_{\text{パイプ}}}{I_{\text{棒}}}=\frac{40^4-30^4}{700^2}=\frac{\num{1750000}}{\num{490000}}\approx3.6 \]
\phtashikame
単位：mm$^4$ どうしの比。\\
大きさ：同じ重さで、パイプは約3.6倍曲がりにくい。材料を中立軸から遠くに置いたからです（H形鋼と同じ考え）。\\
極端な場合：内径を0に近づけると、パイプは棒になり、比は1に近づきます。ただし、うすくしすぎると、パイプは局部的につぶれやすくなります（力-11）。
\kotae 棒の直径 約\qty{26.5}{mm}。パイプの $I$ は棒の約3.6倍
\end{phreidai}
\begin{jibun}
外径\qty{50}{mm}、内径\qty{40}{mm}のパイプと、同じ断面積の丸い棒の $I$ をくらべましょう。
\bkjans{棒の直径は\qty{30}{mm}。$(50^4-40^4)\div30^4\approx4.6$ 倍}
\end{jibun}

\begin{phreidai}[本と置物：重ね合わせ]
\label{ph:力-10:ex10}例題1の本棚の板に、本\qty{30}{kg}を一様に並べ、さらにまん中に\qty{5}{kg}の置物をのせました。まん中のたわみを求めましょう。
\phbamen
等分布の荷重と、まん中の集中荷重が、同時にかかっています。
\phmushi
例題1と同じ（板の重さを無視、たわみは小さい）。
\phzu
⑩の重ね合わせの図。$EI=\pxkq{850.5}{N.m^2}$、$w=\pxkq{367.5}{N/m}$、$F=5\times9.8=\qty{49}{N}$。
\phriyu
たわみの式は1次の式なので、重ね合わせが使えます。等分布だけのたわみと、集中荷重だけのたわみを、別々に求めて足します。どちらもまん中で最大なので、最大どうしを足せます。
\phtateru
\[ \delta=\frac{5wL^4}{384EI}+\frac{FL^3}{48EI} \]
\phtoku
\[ \delta=\qty{2.30}{mm}+\frac{49\times0.8^3}{48\times850.5}=\qty{2.30}{mm}+\qty{0.61}{mm}=\qty{2.92}{mm} \]
\phtashikame
単位：どちらも m（mm）。\\
大きさ：\qty{5}{kg}は本の $\dfrac{1}{6}$ ですが、まん中に集まっているので、たわみは $\dfrac{1}{6}$ より大きい約 $\dfrac{1}{4}$ ふえます。\\
注意：最大の場所がちがう荷重どうしでは、最大どうしを足すと、本当の最大より少し大きめになります（安全がわの見積もり）。
\kotae 約\qty{2.9}{mm}
\end{phreidai}
\begin{jibun}
じぶんでとこう（例題1）の板（$EI=\pxkq{1400}{N.m^2}$、長さ\qty{0.9}{m}、本\qty{40}{kg}を一様に）のまん中に、さらに\qty{8}{kg}の物をのせました。まん中のたわみを求めましょう。
\bkjans{\qty{2.66}{mm}＋\qty{0.85}{mm}で、約\qty{3.5}{mm}}
\end{jibun}

\begin{phreidai}[水の入った配管を、何 m ごとに支える？]
\label{ph:力-10:ex11}水を満たした鋼の配管 SGP 50A（外径\qty{60.5}{mm}、厚さ\qty{3.8}{mm}、$E=\qty{205}{GPa}$）を、間かく\qty{4}{m}の支えにのせます。支えの間のまん中のたわみを求めましょう。間かくを\qty{6}{m}にすると、どうなるでしょう。
\phbamen
工場やビルの天井の配管は、つり金具や台で何 m かおきに支えます。間かくを広げると支えの数がへりますが、配管はたわみます。
\phmushi
支えの所は回れる（単純支持）とし、となりの径間とのつながりは無視します（安全がわ）。弁やフランジの重さは無視し、鋼の密度は\qty{7850}{kg/m^3}、水は\qty{1000}{kg/m^3}とします。
\phzu
外径 $D=\qty{60.5}{mm}$、内径 $d=60.5-2\times3.8=\qty{52.9}{mm}$。1 m あたりの重さ $w$、支えの間かく $L$。
\phriyu
配管の重さと水の重さは、長さにそって一様にかかるので、等分布の単純支持 $\delta=\dfrac{5wL^4}{384EI}$ です。$I$ はパイプの式です。
\phtateru
\[ I=\frac{\pi(0.0605^4-0.0529^4)}{64} \]
\[ w=\Bigl(\frac{\pi}{4}(D^2-d^2)\times7850+\frac{\pi}{4}d^2\times1000\Bigr)\times9.8 \]
\phtoku
\[ I=\qty{2.73e-7}{m^4},\qquad w=(5.31+2.20)\times9.8=\pxkq{73.6}{N/m} \]
\[ L=\qty{4}{m}:\ \delta=\frac{5\times73.6\times4^4}{384\times205\times10^9\times2.73\times10^{-7}}=\qty{4.4e-3}{m} \]
\[ L=\qty{6}{m}:\ \delta=\qty{4.4}{mm}\times\Bigl(\frac{6}{4}\Bigr)^4=\qty{4.4}{mm}\times5.06\approx\qty{22}{mm} \]
\phtashikame
単位：m。\\
大きさ：鋼管の重さ\qty{5.31}{kg/m}は、JIS の表の値と同じです。\qty{4}{m}で\qty{4.4}{mm}は間かくの約 $\dfrac{1}{900}$。\qty{6}{m}では\qty{22}{mm}もたわみ、水がたまる・つなぎ目に力がかかるなどの心配が出ます。\\
極端な場合：間かくを半分の\qty{2}{m}にすると、たわみは $\dfrac{1}{16}$ の\qty{0.27}{mm}です。
\kotae \qty{4}{m}で約\qty{4.4}{mm}、\qty{6}{m}で約\qty{22}{mm}（約5倍）
\end{phreidai}
\begin{jibun}
水を満たした SGP 25A（外径\qty{34.0}{mm}、厚さ\qty{3.2}{mm}、$E=\qty{205}{GPa}$）を、間かく\qty{3}{m}の支えにのせました。まん中のたわみを求めましょう（鋼\qty{7850}{kg/m^3}、水\qty{1000}{kg/m^3}）。
\bkjans{内径\qty{27.6}{mm}、$I=\qty{3.71e-8}{m^4}$、$w\approx\pxkq{29.7}{N/m}$。$\delta\approx\qty{4.1}{mm}$}
\end{jibun}

\begin{phreidai}[長い棚板を、まん中でも支える（不静定のはり）]
\label{ph:力-10:ex12}長さ\qty{1.6}{m}の棚板に、本を1 m あたり\pxkq{367.5}{N/m}で一様に並べます。両はしだけで支えたときと、まん中にも支えを入れたときで、たわみと、まん中の支えが受ける力をくらべましょう（板は例題1と同じで $EI=\pxkq{850.5}{N.m^2}$）。
\phbamen
長い棚の中ほどに、仕切りの板を入れて支える場面です。
\phmushi
板の重さは無視し、支えはどれも板を上へおすだけで、同じ高さにあるとします。
\phzu
\par\vspace{3pt}\noindent\begin{minipage}[c]{0.52\linewidth}
\begin{tikzpicture}[x=1mm, y=1mm]
  \fill[bkOrange!35] (0,0) rectangle (56,2.4); \draw[bkOrange!70!black, line width=0.5pt] (0,0) rectangle (56,2.4);
  \pxkpin{(0,0)} \pxkroller{(28,0)} \pxkroller{(56,0)} \pxkdist{0}{56}{2.4}{$w$}
  \draw[bkGray, line width=0.4pt, <->] (0,-9) -- (28,-9); \node[anchor=north, font=\sffamily\footnotesize] at (14,-9.4) {$L=\qty{0.8}{m}$};
  \draw[bkGray, line width=0.4pt, <->] (28,-9) -- (56,-9); \node[anchor=north, font=\sffamily\footnotesize] at (42,-9.4) {$L=\qty{0.8}{m}$};
\end{tikzpicture}
\end{minipage}\hfill
\begin{minipage}[c]{0.44\linewidth}
支えの間の長さ $L$、まん中の支えの力 $R$、両はしの力 $R_{\mathrm e}$。
\end{minipage}
\phriyu
支えが3つなので、つり合いの式だけでは決まりません（不静定）。⑩の「まん中は下がらない」という条件と重ね合わせを使い、$R=1.25wL$ を使います。
\phtateru
両はしだけ：$\delta=\dfrac{5w(2L)^4}{384EI}$。まん中も支える：$R=1.25wL$、$R_{\mathrm e}=0.375wL$。
\phtoku
\[ \text{両はしだけ：}\ \delta=\frac{5\times367.5\times1.6^4}{384\times850.5}=\qty{3.69e-2}{m}\approx\qty{37}{mm} \]
\[ R=1.25\times367.5\times0.8=\qty{367.5}{N} \]
\[ R_{\mathrm e}=0.375\times367.5\times0.8=\qty{110.25}{N} \]
まん中を支えると、それぞれの径間のたわみは約\qty{0.96}{mm}になります（数値で求めた値。単純支持の\qty{0.8}{m}の板の\qty{2.30}{mm}より小さい）。
\phtashikame
力の和：$367.5+2\times110.25=588=367.5\times1.6$。本の重さの合計と等しい。\\
大きさ：まん中の支えは本の重さの $\dfrac{367.5}{588}=62.5\,\%$ を受けます。「3つで $\dfrac{1}{3}$ ずつ」ではありません。\\
極端な場合：両はしだけでは\qty{37}{mm}もたわみ、本棚として使えません。
\kotae 両はしだけ 約\qty{37}{mm}。まん中の支え\qty{367.5}{N}、両はし\qty{110}{N}ずつ
\end{phreidai}
\begin{jibun}
3つの支えに、長さ\qty{1.0}{m}ずつの2つの径間でのせたはりに、\pxkq{500}{N/m}の等分布荷重がかかっています。まん中の支えと、両はしの支えの力を求めましょう。
\bkjans{まん中\qty{625}{N}、両はし\qty{187.5}{N}ずつ}
\end{jibun}

\begin{phreidai}[両はしを固定したはり：微分方程式を境界条件で解く]
\label{ph:力-10:ex13}長さ $L$ のはりの両はしを、かべにしっかりうめこみました（両端固定）。等分布荷重 $w$ のとき、まん中のたわみが、単純支持の $\dfrac{1}{5}$ になることを示しましょう。
\phbamen
建物のはりは、柱にしっかりつないであり、はしで回ることができません。
\phmushi
はりの重さは $w$ にふくめ、たわみは小さいとします。
\phzu
はしの曲げモーメントを $-M_0$（山の形の向き）とおきます。$x$ の所の曲げモーメントは、単純支持の $M$ に $-M_0$ を加えて $M=\dfrac{w}{2}(Lx-x^2)-M_0$。境界条件は $x=0$ と $x=L$ で $y=0$、$\dfrac{\dd y}{\dd x}=0$。
\phriyu
不静定で、$M_0$ はつり合いだけでは決まりません。たわみの微分方程式 $EI\dfrac{\dd^2y}{\dd x^2}=-M$ を2回積分し、4つの境界条件で $C_1$・$C_2$・$M_0$ を決めます。
\phtateru
\[ EI\frac{\dd y}{\dd x}=-\frac{w}{2}\Bigl(\frac{Lx^2}{2}-\frac{x^3}{3}\Bigr)+M_0x+C_1 \]
\[ EIy=-\frac{w}{2}\Bigl(\frac{Lx^3}{6}-\frac{x^4}{12}\Bigr)+\frac{M_0x^2}{2}+C_1x+C_2 \]
\phtoku
$x=0$ で傾き0・たわみ0より $C_1=0$、$C_2=0$。$x=L$ で傾き0より
\[ -\frac{w}{2}\Bigl(\frac{L^3}{2}-\frac{L^3}{3}\Bigr)+M_0L=0,\qquad M_0=\frac{wL^2}{12} \]
（$x=L$ でたわみ0も、この $M_0$ で成り立ちます。）まん中 $x=\dfrac{L}{2}$ では、
\[ EI\delta=-\frac{w}{2}\Bigl(\frac{L^4}{48}-\frac{L^4}{192}\Bigr)+\frac{wL^2}{12}\cdot\frac{L^2}{8}=-\frac{wL^4}{128}+\frac{wL^4}{96}=\frac{wL^4}{384} \]
\phtashikame
単位：$\dfrac{\pxku{N/m}\times\unit{m^4}}{\pxku{N.m^2}}=\unit{m}$。\\
くらべる：単純支持は $\dfrac{5wL^4}{384EI}$ なので、ちょうど $\dfrac{1}{5}$。\\
数で：本棚の板（$L=\qty{0.8}{m}$）なら\qty{2.30}{mm}が\qty{0.46}{mm}になります。はしの $M_0=\dfrac{wL^2}{12}=\pxkq{19.6}{N.m}$、まん中の $M$ は $\dfrac{wL^2}{8}-\dfrac{wL^2}{12}=\dfrac{wL^2}{24}$。最大の $M$ は、はしに移ります。
\kotae $\delta=\dfrac{wL^4}{384EI}$（単純支持の $\dfrac{1}{5}$）
\end{phreidai}
\begin{jibun}
両端固定のはりのまん中に集中荷重 $F$ をかけると、まん中のたわみは $\dfrac{FL^3}{192EI}$ です。単純支持（$\dfrac{FL^3}{48EI}$）の何分の1でしょう。また、本棚の板（$L=\qty{0.8}{m}$、$EI=\pxkq{850.5}{N.m^2}$）のまん中に\qty{294}{N}をかけたときの、両端固定のたわみを求めましょう。
\bkjans{$48\div192=\frac{1}{4}$。$\frac{FL^3}{192EI}\approx\qty{0.92}{mm}$}
\end{jibun}

%% ================================================================
\bksection{tsumazuki}{まちがえやすいところをチェック}
\begin{tsumazuki}{長さが2倍なら、たわみも2倍？}
\bkng{はりの長さが2倍になると、たわみも2倍になる。}
\bkok{同じ力なら、たわみは長さの3乗で8倍（等分布で1 m あたりが同じなら16倍）。}
うでが長くなって曲げる働きがふえ、曲がる部分も長くなり、角も大きくなるからです。
\end{tsumazuki}
\begin{tsumazuki}{2枚重ねれば8倍かたい？}
\bkng{板を2枚重ねると、高さが2倍なので、たわみは $\dfrac{1}{8}$ になる。}
\bkok{重ねただけでは $\dfrac{1}{2}$。接着してずれないようにすると $\dfrac{1}{8}$。}
重ねただけの板は、重なった面ですべるので、1枚ずつが別々に曲がります。
\end{tsumazuki}
\begin{tsumazuki}{mm$^4$ と m$^4$ の換算}
\bkng{$I=\qty{6750}{mm^4}$ は $\qty{6.75}{m^4}\times10^{-3}$。}
\bkok{$\qty{1}{mm}=\qty{e-3}{m}$ なので $\qty{1}{mm^4}=(10^{-3})^4\,\unit{m^4}=\qty{e-12}{m^4}$。$\qty{6750}{mm^4}=\qty{6.75e-9}{m^4}$。}
4乗の単位は、換算の数も4乗です（2-10 の単位つきの1を4回かける）。
\end{tsumazuki}
\begin{tsumazuki}{2つの $y$}
\bkng{$\sigma=\dfrac{My}{I}$ の $y$ は、はりが下がった量（たわみ）。}
\bkok{$\sigma=\dfrac{My}{I}$ の $y$ は、断面の中で中立軸からはかった距離。}
どちらも「下向きの $y$」ですが、はかり始めがちがいます。たわみの $y$ は、はりの軸が下がった量です。
\end{tsumazuki}
\begin{tsumazuki}{いちばん大きい $M$ は、いつもまん中？}
\bkng{はりの曲げモーメントは、いつもまん中で最大になる。}
\bkok{$M$ が最大になるのは $Q=0$ になる所（荷物の位置・片持ちの根もと・固定したはし）。}
例題3では荷物の所、片持ちはりでは根もとが最大でした。
\end{tsumazuki}

\begin{column}{ガリレオのまちがい}
はりの強さを、はじめて計算で考えたのはガリレオ・ガリレイです。1638年の本『新科学対話』で、かべから出た片持ちの棒が根もとで折れる強さを考えました。でもガリレオは、折れる直前に、断面全体が同じ強さで引っぱられていると考えたので、強さを本当の3倍に見積もってしまいました。断面の上と下でのびちぢみがちがい、中立軸があることに気づいたのは、のちの人たちです。1773年にはクーロンが正しい考え方を示し、1826年のナビエの講義の本で、中立軸が断面の重心を通ることや $\sigma=\dfrac{My}{I}$ の形が、工学で広く使われるようになりました。
\end{column}

%% ================================================================
\Needspace{14\baselineskip}
\bksection{phgear}{工学ではこう使う}
\begin{phkougaku}[建物の床・橋・配管]
\par\vspace{3pt}\noindent\begin{minipage}[c]{0.56\linewidth}
\phboxhead{床のはりは「たわみ」で決まる}
家やビルの床のはりは、こわれない強さ（応力）のほかに、「歩いたときにゆれない」「床が下がって見えない」ことが大切です。日本の建物では、床のはりのたわみをスパンの $\dfrac{1}{250}$〜$\dfrac{1}{300}$ 以下にするのが目安です。木のはりは、長い年月のあいだにたわみがふえる（クリープ）ので、計算のたわみを2倍してくらべます。
\end{minipage}\hfill
\begin{minipage}[c]{0.40\linewidth}
\centering
\begin{tikzpicture}[x=1mm, y=1mm]
  \fill[bkGray!35] (0,0) rectangle (4,26); \fill[bkGray!35] (46,0) rectangle (50,26);
  \draw[bkInk, line width=0.5pt] (0,0) rectangle (4,26) (46,0) rectangle (50,26);
  \fill[bkOrange!40] (4,22) rectangle (46,26); \draw[bkOrange!70!black, line width=0.5pt] (4,22) rectangle (46,26);
  \fill[bkYellow!35] (4,26) rectangle (46,28);
  \bkperson{(18,28)}{12}{bkBlue}
  \bkperson{(31,28)}{11}{bkRed}
  \draw[bkGray, line width=0.4pt, <->] (4,18) -- (46,18); \node[anchor=north, font=\sffamily\footnotesize] at (25,17.5) {スパン};
  \node[anchor=north, font=\sffamily\footnotesize, align=center] at (25,10) {たわみは\\スパンの $\frac{1}{250}$〜$\frac{1}{300}$ 以下};
\end{tikzpicture}
\end{minipage}

\phboxhead{橋の形}
鉄道や道路の橋には、H形鋼より大きい I の形のはり（板を溶接して作る「プレートガーダー」）や、箱の形のはり（箱桁）が使われます。どれも、材料を中立軸から遠い上と下に集めた形です。長い橋では、はりの高さが数 m にもなります。

\phboxhead{配管の支持の間かく}
工場の配管は、例題11のように、たわみと曲げ応力の両方から支持の間かくを決めます。たわみが大きいと、水平のはずの管の中に水や液がたまったり、フランジのつなぎ目に曲げの力がかかって、もれる原因になったりします。保全の点検では、支持金具のゆるみやはずれを見ます。支えが1つはずれると、間かくが2倍になり、たわみは16倍になるからです。

\begin{phketa}[はりの数値]
\phketaline{鋼のヤング率 $E$}{約\qty{205}{GPa}}
\phketaline{すぎのヤング率（繊維の向き、ばらつきが大きい）}{約\qty{7}{GPa}（\qtyrange{6}{12}{GPa}）}
\phketaline{本棚の板のたわみ（例題1）}{約\qty{2}{mm}}
\phketaline{床のはりのたわみの目安}{スパンの $\frac{1}{250}$〜$\frac{1}{300}$}
\phketaline{SGP 50A の重さ（水なし）}{\qty{5.31}{kg/m}}
\phboxhead{答えが桁で合っているか}
たわみが長さの $\dfrac{1}{100}$ をこえたら、計算の単位（mm$^4$ と m$^4$）をまず疑おう。$I$ の換算で $10^{12}$ をまちがえると、答えは何けたもずれます。また、たわみが長さの $\dfrac{1}{10}$ 近くになると「たわみは小さい」という仮定がくずれます。
\end{phketa}
\end{phkougaku}
\phstation{構造力学,材料力学,配管設計,建築構造}{せん断力図（SFD）、曲げモーメント図（BMD）、断面係数、曲げ剛性、たわみ制限、連続ばり}


%% ================================================================
\bkendproblems
\bkproblemgroup{基本}{3}
\bkproblem 両はしで支えた板のまん中に荷物をのせます。次のように変えると、まん中のたわみは何倍になるでしょう。\par
（1）\nobreak 長さを1.5倍\quad（2）\nobreak 厚さを1.5倍\quad（3）\nobreak はばを1.5倍
\bkpans{（1）\nobreak $1.5^3\approx3.4$ 倍（2）\nobreak $1\div1.5^3\approx0.30$ 倍（3）\nobreak $1\div1.5\approx0.67$ 倍}
\bkproblem 長さ\qty{2.4}{m}の板を両はしA・Bで支え、Aから\qty{0.6}{m}の所に\qty{480}{N}の荷物をのせました。A・Bの支える力と、いちばん大きい曲げモーメントを求めましょう。
\bkpans{$R_{\mathrm A}=\qty{360}{N}$、$R_{\mathrm B}=\qty{120}{N}$、$M_{\max}=\pxkq{216}{N.m}$}
\bkproblem はば\qty{40}{mm}、高さ\qty{60}{mm}の長方形の断面の $I$ と $Z$ を、mm の単位で求めましょう。
\bkpans{$I=\qty{7.2e5}{mm^4}$、$Z=\qty{2.4e4}{mm^3}$}
\bkproblem 長さ\qty{0.5}{m}の鋼の平らな棒（はば\qty{6}{mm}、高さ\qty{50}{mm}、立てて使う、$E=\qty{205}{GPa}$）を片持ちにし、先に\qty{300}{N}をかけました。先のたわみと、根もとの曲げモーメント・曲げ応力を求めましょう。
\bkpans{$I=\qty{6.25e-8}{m^4}$、$\delta\approx\qty{0.98}{mm}$、$M=\pxkq{150}{N.m}$、$Z=\qty{2.5e-6}{m^3}$、$\sigma=\qty{60}{MPa}$}
\bkproblem 長さ\qty{2.0}{m}の単純支持のはり（断面 はば\qty{60}{mm}・高さ\qty{90}{mm}）に、\pxkq{1500}{N/m}の等分布荷重がかかっています。いちばん大きいせん断力・曲げモーメント・曲げ応力を求めましょう。
\bkpans{$Q_{\max}=\qty{1500}{N}$、$M_{\max}=\pxkq{750}{N.m}$、$Z=\qty{8.1e-5}{m^3}$、$\sigma\approx\qty{9.3}{MPa}$}
\bkproblem 外径\qty{89.1}{mm}、厚さ\qty{4.2}{mm}の鋼管（SGP 80A）の $I$ を求めましょう。
\bkpans{内径\qty{80.7}{mm}。$I\approx\qty{1.01e6}{mm^4}$}
\bkproblemgroup{チャレンジ}{4}
\bkproblem 長さ\qty{3.0}{m}の鋼の単純支持のはり（$I=\qty{4.0e-6}{m^4}$、$E=\qty{205}{GPa}$）に、\pxkq{2000}{N/m}の等分布荷重と、まん中に\qty{3000}{N}の集中荷重がかかっています。まん中のたわみを求めましょう。
\bkpans{$\qty{2.57}{mm}+\qty{2.06}{mm}\approx\qty{4.6}{mm}$}
\bkproblem 3つの支えに、長さ\qty{1.2}{m}ずつの2つの径間でのせたはりに、\pxkq{800}{N/m}の等分布荷重がかかっています。それぞれの支えの力を求め、力の和が荷重の合計と等しいことを確かめましょう。
\bkpans{まん中\qty{1200}{N}、両はし\qty{360}{N}ずつ。和は\qty{1920}{N}＝$800\times2.4$}
\bkproblem 例題11の SGP 50A（水を満たす、$w=\pxkq{73.6}{N/m}$、$I=\qty{2.73e-7}{m^4}$）のたわみを\qty{2.5}{mm}以下にしたい。支えの間かくは何 m 以下にすればよいでしょう。
\bkpans{$L^4\leqq\frac{384EI\delta}{5w}$ より $L\leqq\qty{3.5}{m}$}
\phnyushi
\bkproblem 長さ\qty{6}{m}、質量\qty{30}{kg}の一様な板を、左はしAと、Aから\qty{4}{m}の点Bで支えました。質量\qty{60}{kg}の人がAから右へ歩きます。（1）人がAから\qty{2}{m}の所にいるとき、A・Bが板を支える力を求めましょう。（2）人がBの外へどこまで進むと、板がAで浮き上がり始めるでしょう。
\bkpans{（1）\nobreak Bのまわりで $4R_{\mathrm A}=294\times1+588\times2$。$R_{\mathrm A}=\qty{367.5}{N}$、$R_{\mathrm B}=\qty{514.5}{N}$（2）\nobreak $294\times1=588x$ より、Bから\qty{0.5}{m}}
\bkproblem 長さ\qty{3.0}{m}、質量\qty{20}{kg}の一様な板を両はしA・Bで支え、Aから\qty{1.0}{m}の所に質量\qty{40}{kg}の荷物をのせました。A・Bが板を支える力を求めましょう。
\bkpans{Bのまわりで $3R_{\mathrm A}=196\times1.5+392\times2.0$。$R_{\mathrm A}\approx\qty{359}{N}$、$R_{\mathrm B}\approx\qty{229}{N}$}

\begin{furikaeri}
\begin{checklist}
\item 板のたわみが、長さ・厚さ・はば・置く物の重さで、どう変わるかを言える。\bkfurref{ph:力-10:sec1}
\item はりを区切って、せん断力・曲げモーメント・たわみを、表計算の足し算で求められる。\bkfurref{ph:力-10:sec4}
\item つり合いと曲げの式から $\sigma=\tfrac{My}{I}$ とたわみの微分方程式を導き、代表的なはりのたわみを計算できる。\bkfurref{ph:力-10:sec8}
\end{checklist}
\end{furikaeri}
\bkfukidashi{sora}{板を立てると100倍曲がりにくいなんて！ 本棚の板の裏に、細い板を立ててはりつけたら、たわまなくなるね。}
